\documentclass[11pt,a4paper]{article}
\usepackage[T1]{fontenc}
\usepackage[utf8]{inputenc}
\usepackage{lmodern,microtype,amsmath,amssymb,amsthm,mathtools}
\usepackage[margin=24mm,headheight=15pt]{geometry}
\usepackage{booktabs,array,longtable,fancyhdr}
\usepackage[hyphens]{url}
\usepackage[hidelinks]{hyperref}
\newtheorem{theorem}{Theorem}[section]
\newtheorem{lemma}[theorem]{Lemma}
\newtheorem{proposition}[theorem]{Proposition}
\newtheorem{corollary}[theorem]{Corollary}
\theoremstyle{remark}\newtheorem{remark}[theorem]{Remark}
\setcounter{tocdepth}{1}
\newcommand{\Li}{\operatorname{Li}_2}
\numberwithin{equation}{section}
\allowdisplaybreaks
\setlength{\emergencystretch}{3em}
\pagestyle{fancy}\fancyhf{}
\fancyhead[L]{\small Completed proofs from Rounds 1--20}
\fancyhead[R]{\small Keio Research Programme}
\fancyfoot[C]{\thepage}
\hypersetup{pdftitle={New Global Closures from Rounds 1--20},pdfauthor={Keio Research Programme}}
\title{New Global Closures from Rounds 1--20\\[5pt]
\large Complete proofs of selected binomial--harmonic identities}
\author{Keio Research Programme\\Synthesis and internal independent verification}
\date{19 November 2026}
\begin{document}
\maketitle
\begin{abstract}
We revisit the incomplete targets in the twenty-round archive and give
complete derivations of 5 original numbered conjectures of Zhi-Wei Sun.
The previously missing steps are evaluated explicitly, including all
infinite-sum boundaries, improper-integral endpoints, and complex logarithm
branches. Two reciprocal-binomial pairs are settled by dilogarithmic reductions,
including an explicit Bloch--Wigner differential calculation. One identity
follows from a harmonic Wilf--Zeilberger pair and its evaluated finite part.
Two further pairs are obtained by differentiating balanced residues of
uniformly convergent hypergeometric accelerations.
All necessary formulas are included in this standalone source. The archive's
six previously completed submissions are excluded from the new-result count.
Other partial advances are not presented as solutions of their original
targets. These are internally checked research proofs; external refereeing
and a comprehensive historical-priority assessment have not been completed.
\end{abstract}
\section*{Scope and reading guide}
\addcontentsline{toc}{section}{Scope and reading guide}
The source numbering throughout is Sun's
\emph{New series involving binomial coefficients (II)},
arXiv:2307.03086v10, revised 6 February 2026~\cite{Sun}.
Conjecture numbers and equation numbers are kept distinct. Each entry below
means the entire original assertion in that numbered conjecture, including
both identities when the source states a pair.
\begin{center}
\begin{tabular}{@{}lllr@{}}\toprule
Original target & Archive origin & Complete scope & Page\\\midrule
Sun 4.21 & Reiwa R04 & Two reciprocal-binomial harmonic identities & \pageref{s421:thm:s421-main}\\
Sun 4.29 & Reiwa R05 & One accelerated harmonic WZ identity & \pageref{thm:sun429-complete}\\
Sun 4.19 & Reiwa R09 & Two reciprocal-binomial harmonic identities & \pageref{s419:thm:sun419-complete}\\
Sun 4.27 & Reiwa R10 & Two accelerated harmonic identities & \pageref{thm:sun427-complete}\\
Sun 4.25 & Reiwa R12 & Two accelerated harmonic identities & \pageref{thm:sun425-complete}\\
\bottomrule\end{tabular}
\end{center}
The proofs use classical beta integrals, dilogarithm identities, exact
telescoping, and elementary estimates. No general-method novelty claim is
made. The proof text, rather than numerical agreement, supplies every
infinite-range conclusion. The final appendix records coverage of the
archive and distinguishes prior completions and an external proof claim.
\clearpage
\section{A complete proof of Sun's Conjecture 4.21}\label{s421:sec:sun421}
Put $H_0=0$, $H_n=\sum_{j=1}^n1/j$, and
\[
 P(k)=77k^2-53k+10,\qquad
 a_k=\frac{1}{k(3k-1)(3k-2)(-2)^k\binom{4k}{k}}.
\]
Define the two absolutely convergent sums
\begin{align}
 S_1&=\sum_{k=1}^{\infty}a_k\left\{
 P(k)(H_{2k-1}-H_{k-1})-\frac{209k^2-113k+18}{2k}\right\},
 \label{s421:eq:S1}\\
 S_2&=\sum_{k=1}^{\infty}a_k\left\{
 P(k)(2H_{4k-1}-H_{2k-1}-H_{k-1})+11k-27+\frac6k\right\}.
 \label{s421:eq:S2}
\end{align}
\begin{theorem}[Sun's Conjecture 4.21]\label{s421:thm:s421-main}
Both series converge absolutely and satisfy the following identities, stated in \cite[p.~20, equations (4.23)--(4.24)]{Sun}:
\begin{equation}\label{s421:eq:original}
 S_1=3\log^2 2\qquad\text{and}\qquad S_2=-\frac{\pi^2}{2}.
\end{equation}
\end{theorem}
The target is not the different Conjecture 4.21 in Sun's older harmonic-series collection. The source here is arXiv:2307.03086v10, revised 6 February 2026.

\subsection{The exact relation between the two sums}
\begin{theorem}\label{s421:thm:sum}
The sums \eqref{s421:eq:S1}--\eqref{s421:eq:S2} satisfy
\begin{equation}\label{s421:eq:sum-result}
 S_1+S_2=3\log^2 2-\frac{\pi^2}{2}.
\end{equation}
It therefore suffices to evaluate $S_1$ separately; this is done below.
\end{theorem}
\begin{proof}
Let $Q(k)=187k^2-59k+6$. Direct addition of the original summands gives
\begin{equation}\label{s421:eq:half-sum}
 \frac{S_1+S_2}{2}
 =\sum_{k\ge1}a_k\left\{P(k)(H_{4k-1}-H_{k-1})-\frac{Q(k)}{4k}\right\}.
\end{equation}
The beta identity
\begin{equation}\label{s421:eq:beta}
 \frac1{k(3k-1)(3k-2)\binom{4k}{k}}
 =\frac{3}{4k(k+1)}B(3k-2,k+2)
\end{equation}
and its parameter derivative imply
\begin{align*}
 &\int_0^1x^{3k-3}(1-x)^{k+1}[-\log(1-x)]\,dx\\
 &\hspace{15mm}=B(3k-2,k+2)(H_{4k-1}-H_{k+1}).
\end{align*}
For fixed $k$, the derivative is dominated near the second parameter by a constant times
$x^{3k-3}(1-x)^{k+1/2}|\log(1-x)|$, an integrable function. The factorial relation \eqref{s421:eq:beta} follows immediately from Euler's beta formula.

Set, on $0<x<1$,
\begin{equation}\label{s421:eq:variables}
 u=-\frac{x^3(1-x)}2,\quad w=\frac{3(1-x)}{4x^3},\quad
 M=\log(1-x),\quad L=\log(1-u),\quad D=\Li(u).
\end{equation}
Here $|u|\le\rho=27/512<1$ and $u<0$, so $L$ and $D=\sum_{n\ge1}u^n/n^2$ are unambiguously real. Partial fractions give
\begin{align*}
 \frac{P(k)}{k(k+1)}&=77+\frac{10}{k}-\frac{140}{k+1},\\
 \frac{P(k)(1/k+1/(k+1))-Q(k)/(4k)}{k(k+1)}
 &=\frac{144}{k+1}-\frac{140}{(k+1)^2}-\frac{147}{4k}+\frac{17}{2k^2}.
\end{align*}
It follows that the half-sum in \eqref{s421:eq:half-sum} equals $\int_0^1I_c(x)\,dx$, where
\begin{align}
 I_c&=w(-AM+B),\label{s421:eq:Ic}\\
 A&=\frac{77u}{1-u}+140+\left(\frac{140}{u}-10\right)L,\notag\\
 B&=-4+\left(\frac{147}{4}-\frac{144}{u}\right)L
       +\left(\frac{17}{2}-\frac{140}{u}\right)D.\notag
\end{align}
For this interchange, both rational coefficients in the two partial fractions are bounded by $1000$. The absolute sum of the integral kernels is bounded by
\[
 \frac{1000w|u|}{1-\rho}(1+|M|)
 =\frac{375}{1-\rho}(1-x)^2(1+|\log(1-x)|),
\]
which is integrable. Thus no formal interchange is being assumed.

Define
\begin{align*}
 V&=\frac{51}{8x}-\frac{51}{16x^2}-\frac{42}{x^5},&
 W&=\frac{15}{2x}-\frac{15}{4x^2}-\frac{42}{x^5},\\
 J&=\frac9x-\frac9{x^2}-\frac{18}{x^5},&
 E&=-\frac{21}{x}+\frac{21}{x^2},&K&=\frac{12}{x}-\frac{12}{x^2},
\end{align*}
and put
\begin{equation}\label{s421:eq:Phi}
 \Phi=(V-\tfrac3{16})D+(W-\tfrac34)ML+JL+EM+K.
\end{equation}
A direct differential certificate is
\begin{equation}\label{s421:eq:cert-combined}
 I_c=\Phi'+\frac{3M}{1+x}.
\end{equation}
For a transparent check, use
\[
 M'=-\frac1{1-x},\quad L'=-\frac{u'}{1-u},\quad
 D'=-L\frac{u'}u,\quad \frac{u'}u=\frac3x-\frac1{1-x}.
\]
The product-rule computation reduces to the following rational identities:
\begin{align*}
 V'&=w(17/2-140/u),\qquad W'=-w(140/u-10),\\
 w(147/4-144/u)+Vu'/u+W/(1-x)
 &=J'-\frac9{16(x-1)}+\frac9{16x},\\
 -w(77u/(1-u)+140)+Wu'/(1-u)
 &=-\frac34L'+\frac3{x+1}+E',\\
 -4w+Ju'/(1-u)+E/(1-x)&=K'.
\end{align*}
They may all be verified by clearing denominators. In the last identity the rational function is $-12/x^2+24/x^3$. These identities establish \eqref{s421:eq:cert-combined} without evaluating an antiderivative by numerical software.

There are no poles in the open interval: $1-u>1$, $1+x>0$, and the coefficients of $\Phi$ only have powers of $x$ in their denominators. At zero, use $D=u+O(x^6)$, $L=-u+O(x^6)$, and $M=-x-x^2/2+O(x^3)$. The five summands of \eqref{s421:eq:Phi} have expansions
\begin{align*}
 (V-\tfrac3{16})D&=21x^{-2}-21x^{-1}+O(x),\\
 (W-\tfrac34)ML&=21x^{-1}-21/2+O(x),\\
 JL&=-9x^{-2}+9x^{-1}+O(x),\\
 EM&=-21x^{-1}+21/2+O(x),\\
 K&=-12x^{-2}+12x^{-1}.
\end{align*}
All divergent and constant terms cancel. Thus $\Phi(0^+)=0$. At one, $D,L=O(1-x)$, $E,K=O(1-x)$, and the only logarithmic factor is $M$. Hence $\Phi(1^-)=0$ as well.

Integrating \eqref{s421:eq:cert-combined} over a compact subinterval and then taking both endpoint limits yields
\[
 \frac{S_1+S_2}{2}=3\int_0^1\frac{\log(1-x)}{1+x}\,dx
 =-3\Li(1/2)=\frac32\log^2 2-\frac{\pi^2}{4}.
\]
The last value follows from
$\Li(z)+\Li(1-z)=\pi^2/6-\log z\log(1-z)$ at $z=1/2$; this reflection identity follows by differentiation and its limit at zero. This proves \eqref{s421:eq:sum-result}.
\end{proof}

\subsection{An exact integral for the first sum}
\begin{theorem}\label{s421:thm:integral}
Define, for $t>0$,
\begin{equation}\label{s421:eq:rR}
 r(t)=\frac{2t^3}{(1+2t^2)^2},\qquad
 R(t)=\frac{27}{4t}+\frac{6t}{1+2t^2}-\frac{12t}{1+4t^2}.
\end{equation}
Then
\begin{equation}\label{s421:eq:remaining}
 S_1=\int_0^\infty\left\{
 R(t)\log(1+r(t)^2)+\frac{12\arctan r(t)}{1+4t^2}
 \right\}\,dt.
\end{equation}
The integral is absolutely convergent and has a positive integrand. The original Conjecture 4.21 is therefore equivalent to the statement that this integral equals $3\log^2 2$.
\end{theorem}
\begin{proof}
We give the complete reduction before evaluating the final integral. Write $h_k=H_{2k-1}-H_{k-1}$. For $-1<u<0$, put $b=\sqrt{-u}$ and define real functions
\[
 N=-\frac{2\arctan b}{b},\qquad
 Z=\frac14\log^2(1-u)+\arctan^2 b.
\]
The endpoint values are $N(0)=-2$, $Z(0)=0$. The basic harmonic generating functions imply
\begin{align}
 \sum_{k\ge1}h_ku^k&=\frac{u(L-N)}{2(1-u)},\label{s421:eq:gen1}\\
 \sum_{k\ge1}\frac{h_k}{k}u^k&=-Z,\label{s421:eq:gen2}\\
 \sum_{k\ge1}\frac{h_k}{k+1}u^k
 &=\frac12-\frac{1+u}{2u}L+N-\frac Zu.\label{s421:eq:gen3}
\end{align}
For example, start from $\sum_{n\ge1}H_nz^n=-\log(1-z)/(1-z)$, take its even part, and use
$h_k=H_{2k}-H_k+1/(2k)$. Integrating the resulting generating function once against $du/u$ gives \eqref{s421:eq:gen2}; integrating against $du$ and dividing by $u$ gives \eqref{s421:eq:gen3}. All series converge absolutely for $|u|<1$. One may use $z=i\sqrt{-u}$ as an intermediate variable: $\log(1-z)$ and $\log(1+z)$ are conjugates whose arguments have positive real parts, so the displayed real formulas have no branch ambiguity.

Apply \eqref{s421:eq:beta} directly to \eqref{s421:eq:S1}, with the variables \eqref{s421:eq:variables}. The additional partial fraction
\[
 \frac{209k^2-113k+18}{2k^2(k+1)}
 =\frac{170}{k+1}-\frac{131}{2k}+\frac9{k^2}
\]
and \eqref{s421:eq:gen1}--\eqref{s421:eq:gen3} show that $S_1=\int_0^1I_1(x)\,dx$, where
\begin{equation}\label{s421:eq:I1}
 \begin{split}
 I_1=w\bigg\{100
  &+\left(\frac92+\frac{240}{u}+\frac{77u}{2(1-u)}\right)L\\
  &+\left(-140-\frac{77u}{2(1-u)}\right)N
    +\left(\frac{140}{u}-10\right)Z-9D\bigg\}.
 \end{split}
\end{equation}
The bound $0<h_k\le1$ and $|u|\le27/512$ justify this termwise integration. For instance, the $k$th integrated absolute term is bounded by $(310/8)(27/512)^{k-1}$ using the beta formula and $0<h_k\le1$.

For use in a second certificate, set
\begin{align*}
 V_1&=\frac{27}{8x^2}-\frac{27}{4x},&
 W_1&=\frac{42}{x^5}+\frac{15}{4x^2}-\frac{15}{2x},\\
 J_1&=\frac{21(1-x)}{x^2},&
 F_1&=-\frac{15}{x}+\frac9{x^2}+\frac{72}{x^5},&
 K_1&=\frac{15}{x}-\frac{15}{x^2},
\end{align*}
and
\[
 \Psi=V_1D+(W_1+\tfrac34)Z+J_1N+F_1L+K_1.
\]
Let
\begin{align*}
 g(x)&=x^3-2x^2+2x-2,\\
 C(x)&=-\frac{6(x^2-3x+1)}{g(x)}+\frac6{x+1},\\
 R_0(x)&=\frac{27}{8}\left(\frac1x+\frac1{1-x}\right)-\frac3{2(x+1)}.
\end{align*}
A second exact differential identity is
\begin{equation}\label{s421:eq:second-cert}
 I_1=\Psi'+C+R_0L+\frac{3N}{2(x+1)}.
\end{equation}
It is checked by the product rule with the derivatives of $L,D$ above and
\[
 N'=-\frac{u'}{u(1-u)}-N\frac{u'}{2u},\qquad
 Z'=\frac{(N-L)u'}{2(1-u)}.
\]
After these substitutions, each coefficient of the formal symbols $D,Z,N,L$ in \eqref{s421:eq:second-cert} is a rational identity; clearing denominators verifies it. These coefficient cancellations are exact polynomial identities after clearing denominators. Notice that $g(x)<0$ on $[0,1]$, since $g(0)=-2$, $g(1)=-1$, and $g'(x)=3x^2-4x+2>0$. Thus no interior pole is introduced.

The boundary limits of $\Psi$ are zero. At zero, use
$D=u+O(x^6)$, $L=-u+O(x^6)$, $Z=-u+O(x^6)$, and
$N=-2-2u/3+O(x^6)$. The five terms are respectively
\[
 O(x),\quad 21x^{-2}-21x^{-1}+O(x),\quad
 -42x^{-2}+42x^{-1}+O(x),
\]
\[
 36x^{-2}-36x^{-1}+O(x),\quad
 -15x^{-2}+15x^{-1}.
\]
They cancel to $O(x)$. At one, $L,D,Z=O(1-x)$, $N$ is bounded, and $J_1,K_1=O(1-x)$; hence $\Psi(1^-)=0$.

Now substitute
\begin{equation}\label{s421:eq:substitution}
 x=\frac{2t^2}{1+2t^2},\qquad
 \theta(t)=\arctan r(t),\qquad\phi(t)=\arctan(2t),\qquad
 \kappa(t)=-\frac3t-3\phi(t).
\end{equation}
Then $\sqrt{-u}=r(t)$, $L=\log(1+r^2)$, and $N=-2\theta/r$. Direct rational differentiation gives
\begin{equation}\label{s421:eq:sub-checks}
 C(x)\frac{dx}{dt}=\frac6t\theta',\qquad
 R_0(x)\frac{dx}{dt}=R(t),\qquad
 \kappa'=\frac{3}{2(x+1)r}\frac{dx}{dt}.
\end{equation}
Therefore the last term of \eqref{s421:eq:second-cert}, after the substitution, is
$-2\theta\kappa'=(-2\kappa\theta)'+2\kappa\theta'$.
The rational part $C(x)dx/dt$ cancels $-6\theta'/t$, leaving
$R(t)L-6\phi\theta'$. Both endpoint limits of $\kappa\theta$ are zero:
$r=O(t^3)$ at zero and $r=O(t^{-1})$ at infinity. Also $\phi\theta$ tends to zero at both ends. A final integration by parts therefore gives
\[
 -6\int_0^\infty\phi\theta'\,dt
 =6\int_0^\infty\phi'\theta\,dt
 =12\int_0^\infty\frac{\theta}{1+4t^2}\,dt,
\]
which proves \eqref{s421:eq:remaining}.

Finally,
\[
 R(t)=\frac{27+138t^2+216t^4}
 {4t(1+2t^2)(1+4t^2)}>0.
\]
The integrand in \eqref{s421:eq:remaining} is positive, is $O(t^3)$ as $t\to0^+$, and is $O(t^{-3})$ as $t\to\infty$. Thus all the improper limits and the final integral converge absolutely.
\end{proof}


\subsection{The missing integral: exact evaluation}
The final evaluation uses a factorization that keeps every logarithm on
a controlled branch. It avoids factoring a cubic into individual roots.
Throughout this subsection, $\operatorname{Log}$ and $\operatorname{Li}_2$
denote their principal branches. For $z$ outside $[1,\infty)$,
\[
 \frac{d}{dz}\operatorname{Li}_2(z)=-\frac{\operatorname{Log}(1-z)}z.
\]
At $z=0$ the apparent singularity is removable.

Set $s=2t$ in \eqref{s421:eq:remaining} and write
\begin{align*}
 r_s&=\frac{s^3}{(s^2+2)^2},&
 u_s&=-i r_s,&
 A_s&=\frac{s^3}{4(s-i)},&
 B_s&=\frac{4(s^2+1)}{(s^2+2)^2},\\
 G_s&=\frac{27}{2s}+\frac{6s}{s^2+2}-\frac6{s-i}.
\end{align*}
For $s>0$, the number $A_s$ lies in the open first quadrant, $B_s>0$,
and $1-u_s=1+i r_s$ also lies in the first quadrant. Consequently
\[
 \operatorname{Log}(1-u_s)
 =\tfrac12\log(1+r_s^2)+i\arctan r_s.
\]
Separating real and imaginary parts of $G_s$ therefore rewrites the
already proved integral exactly as
\begin{equation}\label{s421:eq:s421-complex-integral}
 S_1=\Re\int_0^\infty G_s\operatorname{Log}(1-u_s)\,ds.
\end{equation}
The integral is absolutely convergent: its integrand is $O(s^2)$ at
zero and $O(s^{-2})$ at infinity. Its real part has the stronger
endpoint bounds established in Theorem~\ref{s421:thm:integral}.

The decisive identities are purely rational:
\begin{equation}\label{s421:eq:s421-factorization}
 1-u_s=B_s(1+A_s),\qquad
 G_s=6\frac{A_s'}{A_s}-\frac32\frac{u_s'}{u_s}.
\end{equation}
For the first, multiply through by $(s^2+2)^2$ and use
\[
 (s^2+2)^2+i s^3=(s+i)(s^3+4s-4i).
\]
For the second, use
$A_s'/A_s=3/s-1/(s-i)$ and
$u_s'/u_s=3/s-4s/(s^2+2)$.
Because $B_s$ is positive and $1+A_s$ is in the first quadrant,
there is no integer multiple of $2\pi i$ in the logarithmic identity
\begin{equation}\label{s421:eq:s421-log-split}
 \operatorname{Log}(1-u_s)=\log B_s+\operatorname{Log}(1+A_s).
\end{equation}
The path $u_s$ stays on the negative imaginary axis and tends to zero
at both endpoints. Thus its dilogarithm has the same zero limit at
both endpoints. Integrating \eqref{s421:eq:s421-factorization} and
\eqref{s421:eq:s421-log-split} on a compact subinterval, and then taking
the lower endpoint to zero, gives, for each $R>0$,
\begin{align}
 \Re\int_0^R G_s\operatorname{Log}(1-u_s)\,ds
 &=6I(R^2)-6\Re\operatorname{Li}_2(-A_R)
   +\tfrac32\Re\operatorname{Li}_2(u_R),\label{s421:eq:s421-finite}\\
 I(X)&=\frac12\int_0^X
 \left(\frac3x-\frac1{x+1}\right)
 \left\{\log(1+x)-2\log(1+x/2)\right\}\,dx.
 \label{s421:eq:s421-Ireal}
\end{align}
Indeed $\Re(A_s'/A_s)=3/s-s/(s^2+1)$, and $x=s^2$ produces
\eqref{s421:eq:s421-Ireal}. In \eqref{s421:eq:s421-finite} each term is finite;
the possibly divergent terms will be combined before $R$ tends to
infinity.

An explicit antiderivative, with its lower endpoint already subtracted,
is
\begin{align}
 I(X)={}&-\frac32\operatorname{Li}_2(-X)
       +3\operatorname{Li}_2(-X/2)
       -\operatorname{Li}_2(-X-1)\notag\\
 &-\frac14\log^2(X+1)-\log2\log(X+1)
       +\operatorname{Li}_2(-1).\label{s421:eq:s421-Iclosed}
\end{align}
Differentiation gives the integrand in \eqref{s421:eq:s421-Ireal}; at
$X=0$ the two occurrences of $\operatorname{Li}_2(-1)$ cancel.
For example, the only shifted integral needed is
\[
 \int_0^X\frac{\log(1+x/2)}{x+1}\,dx
 =-\operatorname{Li}_2(-X-1)+\operatorname{Li}_2(-1)
  -\log2\log(X+1).
\]

For clarity, the asymptotic formula used next also follows directly
from a dilogarithm identity. On $\mathbb C\setminus(-\infty,0]$,
\begin{equation}\label{s421:eq:s421-inversion}
 \operatorname{Li}_2(-z)+\operatorname{Li}_2(-1/z)
 =-\frac{\pi^2}{6}-\frac12\operatorname{Log}^2 z.
\end{equation}
Differentiate both sides to see that their difference is constant,
and evaluate at $z=1$, using
$\operatorname{Li}_2(-1)=-\pi^2/12$.
This proves \eqref{s421:eq:s421-inversion} first on the positive axis and
then throughout the indicated connected domain by analytic continuation.
It applies both to the positive real arguments in
\eqref{s421:eq:s421-Iclosed} and to $A_R$ in the first quadrant.

Put $a=\log2$ and $L=\log X$. Equation
\eqref{s421:eq:s421-inversion}, together with
$\operatorname{Li}_2(-1/z)=O(1/|z|)$ as $|z|\to\infty$, gives
\begin{equation}\label{s421:eq:s421-Iasymptotic}
 I(X)=-\frac12 L^2+2aL-\frac32a^2-\frac{\pi^2}{6}+o(1).
\end{equation}
All changes from $\log(X+1)$ to $\log X$ have error
$O((\log X)/X)$. With $X=R^2$,
\[
 A_R=\frac{R^2}{4}\frac1{1-i/R},\qquad
 \operatorname{Log}A_R=L-2a+O(R^{-1}),
\]
so again by \eqref{s421:eq:s421-inversion},
\begin{equation}\label{s421:eq:s421-Aasymptotic}
 -\Re\operatorname{Li}_2(-A_R)
 =\frac12(L-2a)^2+\frac{\pi^2}{6}+o(1).
\end{equation}
The $O(R^{-1})$ logarithmic error, even after multiplication by
$L$, tends to zero. Finally $u_R\to0$, hence
$\operatorname{Li}_2(u_R)\to0$.
Substitute \eqref{s421:eq:s421-Iasymptotic} and
\eqref{s421:eq:s421-Aasymptotic} in \eqref{s421:eq:s421-finite}:
\[
 S_1=6\left(-\frac32a^2+2a^2\right)=3\log^2 2.
\]
The proved relation \eqref{s421:eq:sum-result} now gives
$S_2=-\pi^2/2$. This establishes both assertions of Sun's
Conjecture 4.21.

\paragraph{Verification and provenance.}
The beta identities, the two differential certificates, the changes of
variables, and the two zero boundary values were retained from the
previous round and checked again. The factorization, differential
split, real antiderivative, and dilogarithm asymptotics above supply the
previously missing evaluation. All rational identities were additionally
checked symbolically. Numerical agreement was used only to detect
transcription mistakes; the proof of equality is the exact computation
above. No claim of historical priority is needed for this derivation.

\clearpage
% Fully proved chapter for root's combined manuscript. Requires amsmath,amssymb,amsthm.
% The theorem environments theorem,lemma,proposition and \url should be defined by parent.
\section{A complete evaluation of Sun's Conjecture 4.29}
\label{sec:sun429}

This section proves the single identity stated as Conjecture 4.29,
equation (4.30), in Z.-W. Sun, \emph{New series involving binomial
coefficients (II)}, arXiv:2307.03086v10 (6 February 2026).
Its harmonic-free antecedent and the WZ methodology have substantial
prior history. In particular, the factorial kernel below is related to
the WZ pairs in J. Guillera, \emph{Hypergeometric identities for 10
extended Ramanujan-type series}, Ramanujan J. 15 (2008), 219--234;
arXiv:1104.0396, Identity 4~\cite{Guillera}. We verify all identities and boundary limits
needed here directly. No unevaluated limit from that source is assumed.
This is a complete derivation of the stated target; it is not a claim of
historical priority or of a new general WZ method.

Let $H_0=0$ and $H_m=\sum_{j=1}^m1/j$ for integers $m\ge1$, and put
\[
 P(n)=896n^3+992n^2+296n+21,\qquad
 Q(n)=576n^2+544n+110,
\]
\[
 t_n=\frac{\binom{3n}{n}\binom{4n}{n}^{\!2}}
 {(3n+1)(3n+2)(-2^{14})^n},\qquad
 h_n=4H_{4n}-4H_{2n}+3H_n.
\]
\begin{theorem}\label{thm:sun429-complete}
The series converges absolutely and
\[
 \boxed{\displaystyle
 \sum_{n=0}^{\infty}t_n\bigl(P(n)h_n-Q(n)\bigr)
 =-\frac{256\log2}{\pi}.}
\]
This is the whole assertion of Sun's Conjecture 4.29 in the version
specified above.
\end{theorem}

\subsection{An explicit harmonic WZ pair}
For integers $n,k\ge0$, define
\begin{align}
 B(n,k)&=\frac{(-1)^n(2k)!(2n+2k)!(4n)!}
 {2^{10n+4k}(2n+k)!(n+k)!^2k!n!^2},\label{eq:s429-B}\\
 L(n,k)&=2H_{4n}-3H_{2n}+2H_{2n+2k}-2H_{2k}
       +H_k+H_{n+k}+2H_n,\label{eq:s429-L}\\
 g(n,k)&=\frac{(2n+2k+1)(20n+4k+3)-16kn}{2n+k+1},
 \qquad
 q(n,k)=\frac{28n+12k+12}{2n+k+1},\label{eq:s429-gq}\\
 \mathcal F(n,k)&=B(n,k)\bigl(16nL(n,k)-16\bigr),\label{eq:s429-F}\\
 \mathcal G(n,k)&=B(n,k)\bigl(g(n,k)L(n,k)-q(n,k)\bigr).
 \label{eq:s429-G}
\end{align}
All factorials and harmonic numbers in these definitions have
nonnegative integer indices. All denominators are nonzero.

\begin{lemma}\label{lem:s429-wz}
For every $n,k\ge0$,
\[
 \mathcal G(n,k+1)-\mathcal G(n,k)
 =\mathcal F(n+1,k)-\mathcal F(n,k).
\]
\end{lemma}
\begin{proof}
Here is a rational certificate, including all ingredients required to
check the identity. The factorial ratios are
\begin{align*}
 r_k&:=\frac{B(n,k+1)}{B(n,k)}
 =\frac{(2k+1)(2n+2k+1)}{4(2n+k+1)(n+k+1)},\\
 r_n&:=\frac{B(n+1,k)}{B(n,k)}
 =-\frac{(2n+2k+1)(4n+1)(2n+1)(4n+3)}
 {64(2n+k+1)(2n+k+2)(n+k+1)(n+1)}.
\end{align*}
The harmonic increments reduce to
\begin{align*}
 d_k&:=L(n,k+1)-L(n,k)
 =\frac2{2n+2k+1}+\frac2{n+k+1}-\frac2{2k+1},\\
 d_n&:=L(n+1,k)-L(n,k)\\
 &=\frac2{4n+1}+\frac2{4n+3}-\frac2{2n+1}
   +\frac1{n+1}+\frac2{2n+2k+1}+\frac2{n+k+1}.
\end{align*}
Substitution and clearing denominators give the two rational identities
\begin{align}
 r_kg(n,k+1)-g(n,k)&=16(n+1)r_n-16n,
 \label{eq:s429-cert1}\\
 r_k\{g(n,k+1)d_k-q(n,k+1)\}+q(n,k)
 &-16(n+1)r_nd_n+16r_n-16=0.
 \label{eq:s429-cert2}
\end{align}
For complete transparency, these equalities involve only the rational
functions displayed above; they hold in $\mathbb Q(n,k)$, so the check
is polynomial algebra and has no numerical or finite-range inference.
On expanding the proposed WZ identity after division by $B(n,k)$,
the coefficient of $L(n,k)$ vanishes by
\eqref{eq:s429-cert1}, and its remaining rational part vanishes by
\eqref{eq:s429-cert2}. This proves the lemma.
\end{proof}

\subsection{Absolute bounds and the large-$k$ boundary}
Set
\[
 b_k=B(0,k)=\frac{\binom{2k}{k}^{\!2}}{16^k}
 =\frac{(1/2)_k^2}{k!^2}.
\]
Here $(a)_n=a(a+1)\cdots(a+n-1)$, with $(a)_0=1$.
Factorial duplication gives
\begin{equation}\label{eq:s429-rising}
 \frac{|B(n,k)|}{b_k}
 =\frac{(1/4)_n(1/2)_n(3/4)_n(k+1/2)_n}
 {n!(k+1)_n(k+1)_{2n}}.
\end{equation}
For $0<a\le1$ one has $(a)_n\le n!$, and
$(k+1/2)_n/(k+1)_n\le1$. Also
$\binom{2n}{n}\ge4^n/(2n+1)$, because the largest of the
$2n+1$ binomial coefficients in $(1+1)^{2n}$ is the central one.
Consequently, for all $n,k\ge0$,
\begin{equation}\label{eq:s429-bound1}
 |B(n,k)|\le b_k(2n+1)4^{-n}\le(2n+1)4^{-n}.
\end{equation}
For $n\ge1$, the sharper estimate
\begin{equation}\label{eq:s429-bound2}
 |B(n,k)|\le
 \frac{2b_k(2n+1)4^{-n}}{(k+1)(k+2)}
\end{equation}
follows from
$(k+1)_{2n}\ge(k+1)(k+2)(2n)!/2$.

The elementary bound $H_j\le1+\log j$ for $j\ge1$ implies
\[
 |L(n,k)|\le C\bigl(1+\log(n+k+1)\bigr),\quad
 |g(n,k)|\le56(n+k+1),\quad |q(n,k)|\le28
\]
for an absolute constant $C$. Thus the sums of $\mathcal F(n,k)$ and $\mathcal G(n,k)$ over $n$,
for each fixed $k$, converge absolutely, and their end terms vanish.
The diagonal sums used below also converge absolutely.
In particular, let
\[
 U_k=\sum_{n=0}^{\infty}\mathcal G(n,k).
\]
By \eqref{eq:s429-bound2} and the convergence of polynomially weighted
geometric series,
\begin{equation}\label{eq:s429-rem}
 \sum_{n=1}^{\infty}|\mathcal G(n,k)|
 =O\!\left(\frac{b_k(1+\log(k+1))}{k+1}\right)=o(1).
\end{equation}
The first equality already tends to zero using $b_k\le1$.
At $n=0$ we have
\[
 L(0,k)=2H_k,\qquad q(0,k)=12,\qquad
 g(0,k)=8k+2+\frac1{k+1}.
\]
Stirling's formula gives $b_k=1/(\pi k)+O(k^{-2})$ as $k\to\infty$.
Therefore
\begin{equation}\label{eq:s429-Uk}
 U_k=\frac{16}{\pi}H_k+o(1).
\end{equation}
Summing Lemma~\ref{lem:s429-wz} over $n\ge0$, with the just-proved
absolute convergence and boundary decay, gives
\begin{equation}\label{eq:s429-rec}
 U_{k+1}-U_k=-\mathcal F(0,k)=16b_k.
\end{equation}

\subsection{The exact finite part of the boundary sum}
\begin{lemma}\label{lem:s429-finitepart}
As $K\to\infty$ through positive integers,
\[
 \sum_{k=0}^{K-1}b_k-\frac{H_K}{\pi}
 \longrightarrow\frac{4\log2}{\pi}.
\]
\end{lemma}
\begin{proof}
Stirling's estimate above shows that
$b_k-1/(\pi(k+1))=O(k^{-2})$, so these differences are absolutely
summable. For $0\le z<1$, the binomial series and Euler's beta integral
give
\[
 \sum_{k\ge0}b_kz^k
 =\frac1\pi\int_0^1\frac{dx}{\sqrt{x(1-x)(1-zx)}}.
\]
Tonelli's theorem applies to the nonnegative integrands.
With $z=1-\varepsilon$ and $x=\sin^2\theta$, followed by
$t=\tan\theta$, this becomes
\[
 \sum_{k\ge0}b_k(1-\varepsilon)^k
 =\frac2\pi K_\varepsilon,\qquad
 K_\varepsilon=\int_0^\infty
 \frac{dt}{\sqrt{(1+t^2)(1+\varepsilon t^2)}}.
\]
The substitution $t\mapsto1/(\sqrt\varepsilon\,t)$ preserves this
differential. If $T=\varepsilon^{-1/4}$, then
\[
 K_\varepsilon=2\int_0^T
 \frac{dt}{\sqrt{(1+t^2)(1+\varepsilon t^2)}}.
\]
For $y\ge0$, $0\le1-(1+y)^{-1/2}\le y/2$. Hence
\[
 0\le \operatorname{arsinh}T
 -\int_0^T\frac{dt}{\sqrt{(1+t^2)(1+\varepsilon t^2)}}
 \le\frac\varepsilon2\int_0^T t\,dt
 =\frac{\sqrt\varepsilon}{4}.
\]
Since $\operatorname{arsinh}T=\log(2T)+O(T^{-2})$,
\[
 K_\varepsilon=\log4-\frac12\log\varepsilon
 +O(\sqrt\varepsilon).
\]
On the other hand,
$\sum_{k\ge0}z^k/(k+1)=-\log(1-z)/z$.
Absolute summability of the coefficient differences permits taking
$z\uparrow1$ in their series. The preceding asymptotic therefore yields
\[
 \sum_{k\ge0}\left(b_k-\frac1{\pi(k+1)}\right)
 =\frac{4\log2}{\pi}.
\]
Taking its partial sums is precisely the asserted finite-part limit.
\end{proof}

Combining \eqref{eq:s429-rec}, \eqref{eq:s429-Uk}, and
Lemma~\ref{lem:s429-finitepart} gives the completely evaluated initial
boundary
\begin{equation}\label{eq:s429-U0}
 U_0=\lim_{K\to\infty}
 \left(U_K-16\sum_{k=0}^{K-1}b_k\right)
 =-\frac{64\log2}{\pi}.
\end{equation}

\subsection{Diagonal summation and the original series}
For a finite integer $N\ge1$, summing the WZ identity over the
integer triangular region $0\le k<n<N$ gives the exact identity
\begin{equation}\label{eq:s429-diagonal-finite}
 \sum_{n=0}^{N-1}
 \{\mathcal G(n,n)+\mathcal F(n+1,n)\}
 =\sum_{n=0}^{N-1}\mathcal G(n,0)
   +\sum_{k=0}^{N-1}\mathcal F(N,k).
\end{equation}
Indeed, summing
$\mathcal G(n,n)-\mathcal G(n,0)
=\sum_{k=0}^{n-1}(\mathcal F(n+1,k)-\mathcal F(n,k))$
and adding the missing diagonal $\mathcal F(n+1,n)$ leaves just the
top boundary shown. From \eqref{eq:s429-bound1}, uniformly for $k<N$,
\[
 \sum_{k=0}^{N-1}|\mathcal F(N,k)|
 =O\bigl(N^3(1+\log(N+1))4^{-N}\bigr)\longrightarrow0.
\]
The diagonal series is absolutely convergent by the same bound.
Thus \eqref{eq:s429-diagonal-finite} tends to
\[
 \sum_{n\ge0}\{\mathcal G(n,n)+\mathcal F(n+1,n)\}=U_0.
\]

Finally, factorial cancellation gives
\[
 B(n,n)=(3n+1)(3n+2)t_n,
 \qquad L(n,n)=h_n.
\]
Substituting the displayed rational expressions gives
\begin{align*}
 g(n,n)+16(n+1)r_n|_{k=n}
 &=\frac{P(n)}{4(3n+1)(3n+2)},\\
 -q(n,n)+r_n|_{k=n}\{16(n+1)d_n|_{k=n}-16\}
 &=-\frac{Q(n)}{4(3n+1)(3n+2)}.
\end{align*}
Consequently, at every $n\ge0$, including $n=0$,
\[
 \mathcal G(n,n)+\mathcal F(n+1,n)
 =\frac14t_n\bigl(P(n)h_n-Q(n)\bigr).
\]
Equation \eqref{eq:s429-U0} proves
Theorem~\ref{thm:sun429-complete}. Every boundary term has been
evaluated, and no conjectural hypergeometric special value remains.

\paragraph{Primary sources and verification scope.}
The exact original target is at
\url{https://arxiv.org/html/2307.03086v10}, Conjecture 4.29.
The classical WZ antecedent is available at
\url{https://arxiv.org/pdf/1104.0396}.
The rational identities in this proof were additionally checked by
independent symbolic simplification; finite arithmetic checks are only
transcription checks and are not the proof of an infinite identity.

\clearpage
\section{Both identities in Sun's Conjecture 4.19}\label{s419:sec:sun419-complete}
Let $H_0=0$, $H_n=\sum_{j=1}^n1/j$, and
\[
 K=L(2,\chi_{-3})=\sum_{j\ge0}\left\{\frac1{(3j+1)^2}-\frac1{(3j+2)^2}\right\},
 \qquad P(k)=(k-1)(13k-2).
\]
Define
\begin{align}
 a_k&=\frac{9^{k-1}}{k(3k-1)(3k-2)\binom{4k}{k}},\label{s419:ak}\\
 S_1&=\sum_{k\ge1}a_k\{P(k)(3H_{2k-1}-4H_{k-1})+12(4k-1)\},\label{s419:S1}\\
 S_2&=\sum_{k\ge1}a_k\{P(k)(8H_{4k-1}-9H_{2k-1})-160k+84\}.\label{s419:S2}
\end{align}
The exact source is Sun \cite[v10, p.~20, Conjecture 4.19, (4.21)--(4.22)]{Sun}.
The conjecture is dated 2025-02-06. We prove its complete assertion:
\begin{theorem}[Sun's Conjecture 4.19]\label{s419:thm:sun419-complete}
\begin{equation}\label{s419:conj}
 S_1=\frac{27}4K,\qquad
 S_2=\frac{4\pi\log3}{\sqrt3}-\frac{93}4K.
\end{equation}
\end{theorem}
There is no omitted congruence in this numbered conjecture. The distinct
preceding Conjecture 4.17 is not assumed. In particular its harmonic-free
series value is not used as an unproved lemma below.

For $k\ge1$ the exact ratio is
\[
 \frac{a_{k+1}}{a_k}
 =\frac{27k(3k-1)(3k-2)}{4(4k+1)(4k+2)(4k+3)}<q:=\frac{243}{256}<1,
 \qquad a_1=\frac18.
\]
Thus $a_k\le q^{k-1}/8$. The bound $H_m\le m$ proves absolute convergence
of both original series, even with crude cubic polynomial majorants.

\subsection{A fully evaluated linear combination}
\begin{theorem}\label{s419:relation}
Without assuming either equality in \eqref{s419:conj},
\begin{equation}\label{s419:Tvalue}
 T:=3S_1+S_2=\frac{4\pi\log3}{\sqrt3}-3K.
\end{equation}
This reduces the remaining work to a separate evaluation of $S_1$, which is completed below.
\end{theorem}

\subsubsection{Beta conversion and generating functions}
For $k\ge1$, the factorial identity
\begin{equation}\label{s419:beta}
 \frac1{k(3k-1)(3k-2)\binom{4k}{k}}
 =\frac3{4k(k+1)}B(3k-2,k+2)
\end{equation}
uses $B(a,b)=\int_0^1x^{a-1}(1-x)^{b-1}\,dx$.
Differentiating the beta integral in its second parameter gives
\begin{align}
 &\int_0^1x^{3k-3}(1-x)^{k+1}[-\log(1-x)]\,dx\notag\\
 &\hspace{20mm}=B(3k-2,k+2)(H_{4k-1}-H_{k+1}).\label{s419:betader}
\end{align}
For parameter perturbations of modulus less than $1/2$, domination by a
constant times $x^{3k-3}(1-x)^{k+1/2}|\log(1-x)|$ justifies this derivative.
For $0<x<1$ put
\begin{equation}\label{s419:vars}
 u=9x^3(1-x),\quad W=\frac{1-x}{12x^3},\quad
 M=\log(1-x),\quad L=\log(1-u),\quad D=\Li(u).
\end{equation}
Here $0<u\le q<1$, and $\Li(u)=\sum_{k\ge1}u^k/k^2$. All logarithms and
the dilogarithm have unambiguous real values. Absolute summability follows
from $a_k\le q^{k-1}/8$ and polynomial bounds on the harmonic numbers.
This also justifies the interchange of each integrated beta series,
including its logarithmic derivative, with summation. Equivalently, the
power series below converge on the compact range $0\le u\le q$; their
apparently singular expressions at $u=0$ are removable.

Adding the original summands gives
\[
 T=\sum_{k\ge1}a_k\{P(k)(8H_{4k-1}-12H_{k-1})-16k+48\}.
\]
Use $H_{k+1}-H_{k-1}=1/k+1/(k+1)$ in \eqref{s419:betader}. The relevant
coefficient decompositions are
\begin{align*}
 \frac{P(k)}{k(k+1)}&=13+\frac2k-\frac{30}{k+1},\\
 \frac{P(k)}{k(k+1)}\left(\frac{12}k+\frac8{k+1}\right)
 +\frac{-16k+48}{k(k+1)}
 &=\frac{384}{k+1}-\frac{240}{(k+1)^2}-\frac{140}k+\frac{24}{k^2}.
\end{align*}
The elementary generating functions
\[
 \sum_{k\ge1}H_ku^k=-\frac{L}{1-u},\quad
 \sum_{k\ge1}\frac{H_k}ku^k=D+\frac{L^2}2,\quad
 \sum_{k\ge1}\frac{H_k}{k+1}u^k=\frac{L^2}{2u}
\]
now give
\begin{equation}\label{s419:Tintegral}
 T=\int_0^1 I(x)\,dx,
\end{equation}
where, with $A=13u/(1-u)-2L+30L/u+30$,
\begin{equation}\label{s419:I}
 I=W\left\{-8MA+4\left(1-\frac{15}u\right)(4D-L^2)
       +\left(\frac{52}{1-u}+140-\frac{384}u\right)L-144\right\}.
\end{equation}
These equations may also be checked by comparing every coefficient of
$u^k$ before integration.

\subsubsection{An explicit boundary certificate}
Set
\begin{align*}
 R(x)&=\frac1{24x^2}-\frac1{12x}-\frac1{36x^5},&
 U(x)&=-\frac{4(x-1)}{x^2},&
 V(x)&=\frac{4(1-u)}{9x^5},\\
 C(x)&=R(4L^2-16D-16ML)+UM+VL,&
 h(x)&=\frac{3(3x^2-6x+2)}{1-u}.
\end{align*}
The following differential identity is exact on the whole open interval:
\begin{equation}\label{s419:certificate}
 I(x)=C'(x)+h(x)(2M-L).
\end{equation}
One verifies it by using $M'=-1/(1-x)$, $L'=-u'/(1-u)$ and
$D'=-Lu'/u$, and equating coefficients of $M,L,D,M L,L^2$.
It is a rational identity after these differentiation rules are imposed;
clearing denominators gives an identically zero numerator. These exact cancellations were independently rechecked in the present pass.

The boundary terms, not just their formal derivatives, matter. As $x\to0^+$,
the five components of $C$ respectively have the expansions
\begin{align*}
 4RL^2&=O(x),&-16RD&=4x^{-2}-4x^{-1}+O(x),\\
 -16RML&=4x^{-1}-2+O(x),&UM&=-4x^{-1}+2+O(x),\\
 VL&=-4x^{-2}+4x^{-1}+O(x).
\end{align*}
They cancel, so $C(x)\to0$. At $x=1-\varepsilon$, we have
$u,L,D=O(\varepsilon)$, $M=\log\varepsilon$, $U=O(\varepsilon)$ and
bounded $R,V$, giving $C(x)\to0$ there too.
The residual integrand is integrable at both endpoints. Hence
\begin{equation}\label{s419:reducedT}
 T=\int_0^1\frac{3(3x^2-6x+2)}{1-9x^3(1-x)}
       \log\frac{(1-x)^2}{1-9x^3(1-x)}\,dx.
\end{equation}

\subsubsection{Evaluation, including the nonmonotone substitution}
Let
\[
 Y=\frac{6x^2-3x-1}{1-x}.
\]
Direct algebra gives
\[
 1-u=\frac{(1-x)^2}{4}(Y^2+3),\qquad
 h(x)=-\frac{6Y'}{Y^2+3},\qquad Y(0)=-1,\quad Y(1^-)=+\infty.
\]
This map is \emph{not} monotone: $Y'$ vanishes at $1-1/\sqrt3$.
There is nevertheless no change-of-variables gap. Apply the chain rule to
an antiderivative of the continuous function
$\log((y^2+3)/4)/(y^2+3)$; the backtracked portion cancels. We obtain
\begin{align}
 T&=6\int_{-1}^\infty\frac{\log((y^2+3)/4)}{y^2+3}\,dy\notag\\
 &=2\sqrt3\int_{-\pi/6}^{\pi/2}
       \{\log(3/4)-2\log\cos\theta\}\,d\theta.\label{s419:theta}
\end{align}
To evaluate this, define $\operatorname{Cl}_2(\phi)=\sum_{n\ge1}\sin(n\phi)/n^2$.
Integrating the Abel-regularized Fourier series of $\log(2\cos\theta)$,
then taking its integrable endpoint limit, gives
\[
 \int_0^a\log\cos\theta\,d\theta
 =-a\log2+\tfrac12\operatorname{Cl}_2(\pi-2a),\quad 0\le a\le\pi/2.
\]
The absolutely convergent sine series gives
$\operatorname{Cl}_2(2\pi/3)=\sqrt3 K/2$. Thus
\[
 \int_{-\pi/6}^{\pi/2}\log\cos\theta\,d\theta
 =-\frac{2\pi}3\log2+\frac{\sqrt3}4K.
\]
Substitution in \eqref{s419:theta} proves Theorem~\ref{s419:relation}.

\subsection{An exact integral for the first sum}
The following reduction prepares the integral evaluated in the next subsection.
For $t>0$ put
\begin{align*}
 v(t)&=\frac{3t^3}{(1+t^2)^2},&
 \ell_-(t)&=\log(1-v(t)),&\ell_+(t)&=\log(1+v(t)),\\
 p_-(t)&=t^4-3t^3+2t^2+1,&
 p_+(t)&=t^4+3t^3+2t^2+1,\\
 h_-(t)&=\frac{t^2-2t+2}{p_-(t)},&
 h_+(t)&=\frac{t^2+2t+2}{p_+(t)}.
\end{align*}
Since $v^2\le243/256<1$, both quartics are positive on this domain and all
logarithms are real.
\begin{proposition}\label{s419:remaining}
The first original sum has the exact representation
\begin{equation}\label{s419:remainingI}
 S_1=\frac34\int_0^\infty
 \{h_-(\ell_-+4\ell_+)-h_+(4\ell_-+\ell_+)\}\,dt.
\end{equation}
The integrand is $O(t^3)$ at zero and $O(t^{-3})$ at infinity.
The entire original conjecture is therefore equivalent to the assertion
that the right side of \eqref{s419:remainingI} equals $27K/4$.
\end{proposition}
\begin{proof}
In the beta integral use $x=t^2/(1+t^2)$, so that
$u=v^2$ and $W\,dx=dt/(6t^5)$. First rewrite the harmonic coefficient as
\[
 3H_{2k-1}-4H_{k-1}=3H_{2k}-4H_k+\frac5{2k}.
\]
Let $A_1(z)=\sum H_kz^k$, $B_1(z)=\sum H_kz^k/k$, and
$C_1(z)=\sum H_kz^k/(k+1)$, with the explicit formulas in the preceding subsection.
Even-index extraction gives, for $u=v^2$,
\begin{align*}
 A_2(u):=\sum H_{2k}u^k&=\tfrac12(A_1(v)+A_1(-v)),\\
 B_2(u):=\sum H_{2k}u^k/k&=\tfrac12\Li(u)+\tfrac12(\ell_-^2+\ell_+^2),\\
 C_2(u):=\sum H_{2k}u^k/(k+1)
 &=\frac{-(1-v)\ell_--(1+v)\ell_++(\ell_-^2+\ell_+^2)/2}{u}.
\end{align*}
The correction term is controlled by
\[
 \frac{5P(k)}{2k^2(k+1)}+\frac{12(4k-1)}{k(k+1)}
 =-\frac{109}{2k}+\frac5{k^2}+\frac{135}{k+1}.
\]
In the resulting generating function the dilogarithms cancel exactly.
Consequently $S_1=\int_0^\infty J(t)\,dt$, where
\begin{align*}
 r(t)&=-\frac1{36t^{10}}-\frac5{36t^8}-\frac5{18t^6}
       -\frac{17}{72t^4}-\frac5{36t^2}-\frac5{72},\\
 c_-(t)&=\frac1{6t^5}\left\{-\frac{39}{2(1-v)}+\frac{52}{1-v^2}
       +\frac{90(1-v)}{v^2}+\frac{109}2-\frac{135}{v^2}\right\},\\
 c_+(t)&=\frac1{6t^5}\left\{-\frac{39}{2(1+v)}+\frac{52}{1-v^2}
       +\frac{90(1+v)}{v^2}+\frac{109}2-\frac{135}{v^2}\right\},\\
 J(t)&=r'(t)(\ell_-^2+\ell_+^2+8\ell_-\ell_+)
           +c_-\ell_-+c_+\ell_+-\frac{45}{2t^5}.
\end{align*}
Here $r'=(-1+15/v^2)/(6t^5)$, an identity that also checks the Jacobian
and the overall factor $9^{k-1}$.

For an explicit second boundary certificate set
\begin{align*}
 e(t)&=-\frac5{6t^2}-\frac{11}{12t^4}+\frac5{6t^6}
          +\frac5{12t^8}+\frac1{12t^{10}},\\
 o(t)&=\frac1t-\frac3{2t^3}-\frac3{2t^5}-\frac1{2t^7},\\
 E(t)&=r(t)(\ell_-^2+\ell_+^2+8\ell_-\ell_+)
       +(e+\tfrac7{12})(\ell_-+\ell_+)+o(\ell_+-\ell_-)
       +\frac{9(t^2+1)}{4t^4}.
\end{align*}
Direct differentiation gives the rational-coefficient identity
\begin{equation}\label{s419:secondcert}
 J(t)=E'(t)+\tfrac34\{h_-(\ell_-+4\ell_+)-h_+(4\ell_-+\ell_+)\}.
\end{equation}
For verification use $\ell_-'=-v'/(1-v)$ and $\ell_+'=v'/(1+v)$;
each coefficient of $\ell_-^2,\ell_+^2,\ell_-\ell_+,\ell_-,\ell_+,1$
then agrees rationally. At zero, expand
$v=3t^3-6t^5+9t^7-12t^9+O(t^{11})$ and the two logarithms through
degree ten. All negative powers and the constant term of $E$ cancel,
leaving $E=O(t)$. At infinity $v=O(t^{-1})$, so $E=O(t^{-2})$.
Thus both boundary limits vanish. Integrating \eqref{s419:secondcert} proves
\eqref{s419:remainingI}. Its stated convergence estimates follow by expanding
$v,\ell_\pm,h_\pm$ at the two endpoints.
\end{proof}


% New missing-step proof for Sun 4.19. The parent document supplies
% the archived beta reductions and the proved relation 3 S_1+S_2.
% Uses amsmath, amssymb, amsthm, booktabs. No external figures/data needed.
\subsection{Evaluation of the remaining integral by Bloch--Wigner differentials}
\label{s419:sec:sun419-new-evaluation}
We now evaluate the integral in the preceding reduction. Write
\[
 \omega=\frac{-1+i\sqrt3}{2},\qquad
 \mathcal D(z)=\Im\operatorname{Li}_2(z)
                 +\arg(1-z)\log|z|.
\]
Here $\mathcal D$ is the single-valued continuous Bloch--Wigner
dilogarithm, with $\mathcal D(0)=\mathcal D(1)=\mathcal D(\infty)=0$.
In particular its value does not depend on choosing a continuous complex
logarithm along any of the paths below. For nonzero differentiable
complex functions define the real differential
\[
 \eta(f,g)=\log|f|\,\Im\frac{dg}{g}
              -\log|g|\,\Im\frac{df}{f}.
\]
It is alternating, additive in each argument under multiplication,
and unchanged when an argument is multiplied by a constant of modulus
one. Direct differentiation of the defining dilogarithm gives
\begin{equation}\label{s419:eq:sun419-bw-differential}
 d\mathcal D(z)=\eta(z,1-z).
\end{equation}
We will use the resulting elementary identity
\begin{equation}\label{s419:eq:sun419-eta-difference}
 \eta(f,g)=\eta(f,g-f)-\eta(g,g-f)-d\mathcal D(f/g).
\end{equation}
These identities are applied away from the isolated zeros or poles and
then integrated by taking one-sided limits. All singularities introduced
below are logarithmic, or have a continuous dilogarithm limit, so this
procedure requires no principal-value convention.

Put
\[
 q=t^2+(\omega-1)t+\bar\omega,\quad p=q(-t),\quad
 Q=\bar q,\quad P=\bar p,\quad d=1+t^2.
\]
Complex conjugation here means conjugating coefficients, with $t$ real.
Direct multiplication gives
\[
 |q|^2=p_-(t),\quad |p|^2=p_+(t),\quad
 \Im\frac{q'}q=-\frac{\sqrt3}{2}h_-(t),\quad
 \Im\frac{p'}p=\frac{\sqrt3}{2}h_+(t).
\]
The positive quartics in the preceding reduction show that these
quadratics have no zeros on the real axis. Consequently, if
\begin{equation}\label{s419:eq:sun419-Tform}
 \mathcal T=\tfrac12\eta(Q/d,q/d)+\tfrac12\eta(P/d,p/d)
                          +4\eta(P/d,q/d),
\end{equation}
then the previously established integral is exactly
\begin{equation}\label{s419:eq:sun419-S1-T}
 S_1=-\sqrt3\int_0^\infty\mathcal T.
\end{equation}
Indeed $\log|q/d|=\ell_-/2$ and
$\log|p/d|=\ell_+/2$, which verify every coefficient in this formula.

The useful differences of quadratics are linear:
\[
 q-Q=i\sqrt3(t-1),\quad p-P=-i\sqrt3(t+1),\quad
 q-P=-3t-i\sqrt3,\quad p-Q=3t-i\sqrt3.
\]
Applying \eqref{s419:eq:sun419-eta-difference}, using conjugation, and
collecting the terms containing $d$, therefore yields
\begin{equation}\label{s419:eq:sun419-first-reduction}
 \mathcal T=d\mathcal B+\eta(q,F(t))+\eta(p,F(-t)),
\end{equation}
where
\[
 \mathcal B=-\tfrac12\mathcal D(Q/q)-\tfrac12\mathcal D(P/p)
                         -4\mathcal D(P/q),\qquad
 F(t)=\frac{(1+t^2)^5}
 {i\sqrt3(t-1)(3t+i\sqrt3)^4}.
\]
For example the two self terms leave
$-\eta(q,i\sqrt3(t-1))-\eta(p,-i\sqrt3(t+1))$;
the cross term leaves
$-4\eta(q,3t+i\sqrt3)-4\eta(p,3t-i\sqrt3)$;
and the denominators $d$ contribute
$5\eta(q,d)+5\eta(p,d)$. This proves the displayed collection with its
constants as well as its exponents.

\paragraph{Four elementary identities remove the quadratic factors.}
Let $a=-i/\sqrt3$, $b=\bar\omega$, $c=-\omega$, and set
\begin{align*}
 v_1&=t-b,&
 v_2&=\tfrac{3+i\sqrt3}{2}t,&
 v_3&=3(t+\omega),\\
 v_4&=\tfrac{9+3i\sqrt3}{2},&
 v_5&=\tfrac{9+9i\sqrt3}{2}(t-1),&
 l&=t-a.
\end{align*}
The following four polynomial identities can each be checked by
multiplying two linear factors:
\begin{equation}\label{s419:eq:sun419-four-polynomials}
 v_j(t-b)-v_{j+1}l=c_jq\quad(1\le j\le4),
\end{equation}
where
\[
 (c_1,c_2,c_3,c_4)=
 \left(\bar\omega,\frac{-3+i\sqrt3}{2},3,
                         -\frac{9+9i\sqrt3}{2}\right).
\]
Define rational functions
\[
 z_j=\frac{v_j(t-b)}{v_{j+1}l},\qquad
 g_j=\frac{-c_j}{v_{j+1}l};
 \quad\text{thus}\quad 1-z_j=qg_j.
\]
Writing $\gamma=1-\omega$, we also have
\[
 d-q=\gamma(t-b),\qquad
 F(t)=-\left(\frac{d}{\gamma(t-b)}\right)^5
                         \prod_{j=1}^4z_j.
\]
The constant in the latter equality follows, for example, from
$\gamma^5(9+9i\sqrt3)/(2i\sqrt3\,81)=-1$.
By bilinearity and \eqref{s419:eq:sun419-bw-differential},
\begin{equation}\label{s419:eq:sun419-q-elimination}
 \eta(q,F)=d\mathcal A(t)+\mathcal L(t),
\end{equation}
where
\begin{align*}
 \mathcal A(t)&=-5\mathcal D(q/d)-\sum_{j=1}^4\mathcal D(z_j),\\
 \mathcal L(t)&=-5\eta(d,\gamma(t-b))+
                                      \sum_{j=1}^4\eta(z_j,g_j).
\end{align*}
For completeness,
$d\mathcal D(q/d)=-\eta(q,d/(\gamma(t-b)))-\eta(d,\gamma(t-b))$,
whereas $d\mathcal D(z_j)=-\eta(q,z_j)+\eta(z_j,g_j)$;
these two identities prove \eqref{s419:eq:sun419-q-elimination} directly.
Let $\iota(t)=-t$. Its pullback $\iota^*\mathcal L$ means that every
rational function in $\mathcal L$ is evaluated at $-t$ and then
differentiated with respect to the positive variable $t$.
Reflection gives the identical formula for $\eta(p,F(-t))$, with
$\mathcal L$ replaced by $\iota^*\mathcal L$.

\paragraph{Dilogarithm boundary values.}
Write $D_\omega=\mathcal D(\omega)$,
$\rho=(1+i\sqrt3)/2=-\bar\omega$, and
$\zeta=(1+i/\sqrt3)/2=1/(1-\omega)$.
The standard elementary symmetries, which also follow by differentiating
\eqref{s419:eq:sun419-bw-differential} and evaluating at a real point, are
\[
 \mathcal D(\bar z)=-\mathcal D(z),\quad
 \mathcal D(1-z)=-\mathcal D(z),\quad
 \mathcal D(1/z)=-\mathcal D(z),\quad
 \mathcal D(z)+\mathcal D(-z)=\tfrac12\mathcal D(z^2).
\]
They give $\mathcal D(\rho)=3D_\omega/2$ and
$\mathcal D(\zeta)=D_\omega$.
At $t=0$, each of $Q/q,P/p,P/q$ equals $\bar\omega$;
at infinity each tends to $1$. Hence
\[
 \mathcal B(\infty)-\mathcal B(0)=-5D_\omega.
\]
Both $q/d$ and $p/d$ go from $\bar\omega$ to $1$.
The four $z_j$ have the following endpoint limits in the Riemann sphere:
\[
\begin{array}{c|cccc}
 &z_1&z_2&z_3&z_4\\\hline
 t\to0&\infty&0&\rho&\omega\\
 t\to+\infty&\bar\zeta&\zeta&\infty&0\\
 t\to-\infty&\bar\zeta&\zeta&\infty&0
\end{array}
\]
It follows that $\mathcal A(+\infty)-\mathcal A(0)
=\mathcal A(-\infty)-\mathcal A(0)=-5D_\omega/2$.
Thus the total boundary contribution in
\eqref{s419:eq:sun419-first-reduction}--\eqref{s419:eq:sun419-q-elimination}
is $-10D_\omega$.

\paragraph{The remaining linear differentials.}
Here we give explicit values so that no unevaluated dilogarithm
combination remains. Write $E_\alpha=t-\alpha$ and
\[
 J(\alpha,\beta)=\int_0^\infty
 \{\eta(E_\alpha,E_\beta)+\eta(E_{-\alpha},E_{-\beta})\},
 \qquad
 J_C(\alpha)=\int_0^\infty
 \{\eta(3,E_\alpha)+\eta(3,E_{-\alpha})\}.
\]
Multiplication by the constant $-1$ under reflection does not change
$\eta$, so these are precisely the paired integrals needed for
$\mathcal L(t)$ and its reflection.
To evaluate them, define
\[
 A(\alpha)=
 \begin{cases}\arg(-\alpha),&\Im\alpha\ne0,\\0,&\alpha\in\mathbb R.
 \end{cases}
 \qquad \sigma(\alpha)=A(\alpha)+A(-\alpha).
\]
The arguments in this definition are principal ones. A linear factor
with nonreal root remains in a fixed open half-plane along the real
integration path, so its total change of argument is $-A(\alpha)$.
For real roots its argument derivative is identically zero away from
the root; this explains the second case and excludes any spurious
argument jump at $t=1$.
Applying \eqref{s419:eq:sun419-bw-differential} to
$(t-\alpha)/(\beta-\alpha)$ gives, for distinct roots,
\begin{align}
 J(\alpha,\beta)
 &=-2\mathcal D\left(\frac{\alpha}{\alpha-\beta}\right)
 +\log|\beta-\alpha|\{\sigma(\alpha)-\sigma(\beta)\},
 \label{s419:eq:sun419-J}\\
 J_C(\alpha)&=-\log3\,\sigma(\alpha).
 \label{s419:eq:sun419-JC}
\end{align}
These improper integrals converge; a real root produces at most a
logarithmic singularity. The dilogarithm is continuous at any crossing
of $0$ or $1$, so no extra boundary term arises there.

For clarity, expansion of the products in $\mathcal L$ gives the
following complete linear differential. The constant in $\eta(3,E)$
is the positive constant function $3$:
\begin{align*}
 \mathcal L={}&\tfrac52\eta(3,E_i)+5\eta(E_b,E_i)
 +\tfrac52\eta(3,E_{-i})+5\eta(E_b,E_{-i})
 +2\eta(3,E_b)-2\eta(E_b,E_0)\\
 &+5\eta(E_a,E_b)+\tfrac12\eta(3,E_0)
 +\tfrac12\eta(3,E_a)+\eta(E_a,E_0)
 +\tfrac12\eta(3,E_c)+\eta(E_c,E_0)\\
 &-\eta(E_b,E_c)+\eta(E_a,E_c)
 +\tfrac12\eta(3,E_1)-\eta(E_b,E_1).
\end{align*}
This expansion uses only
$|\gamma|=|{(3+i\sqrt3)}/{2}|=\sqrt3$,
$|v_4|=3\sqrt3$, and
$|{(9+9i\sqrt3)}/{2}|=9$.
All paired integrals in this expansion vanish except those in the next
table; the assertions follow directly from
\eqref{s419:eq:sun419-J}--\eqref{s419:eq:sun419-JC}:
\[
\begin{array}{c|c}
 \text{paired integral}&\text{value}\\\hline
 J(b,i)+J(b,-i)&-\tfrac32D_\omega\\
 J(a,b)&2D_\omega-\pi\log3/6\\
 J(b,c)&-3D_\omega\\
 J(a,c)&-2D_\omega+\pi\log3/6\\
 J(b,1)&-2D_\omega-\pi\log3/6\\
 J_C(b)&\pi\log3/3\\
 J_C(c)&-\pi\log3/3
\end{array}
\]
For the first row the two logarithmic contributions cancel, and
\[
 \mathcal D\left(\frac{b}{b-i}\right)+
 \mathcal D\left(\frac{b}{b+i}\right)
 =-\mathcal D(-ib)-\mathcal D(ib)
 =-\tfrac12\mathcal D(-\omega)=\tfrac34D_\omega.
\]
The other nonzero dilogarithm arguments in
\eqref{s419:eq:sun419-J} are, respectively,
$\bar\omega,\rho,\omega,\zeta$.
The argument sums are $\sigma(a)=\sigma(i)=\sigma(-i)=0$,
$\sigma(b)=-\pi/3$, $\sigma(c)=\pi/3$,
and $\sigma(0)=\sigma(1)=0$, which also verify every logarithmic term.
Multiplying the table rows by the coefficients in $\mathcal L$ gives
\[
 \int_0^\infty(\mathcal L+\iota^*\mathcal L)
 =\tfrac{11}{2}D_\omega.
\]
Both differentials are integrated in the positive $t$ direction.
The $\pi\log3$ terms cancel
exactly: the linear-pair terms contribute $-\pi\log3/2$ and the
constant-pair terms contribute $+\pi\log3/2$.

\paragraph{Conclusion and interior limits.}
The rational functions introduced above can have real poles at
$t=0$ or $t=1$. At $t=1$ only the explicitly displayed factor $t-1$
vanishes. Each $\mathcal D$ term remains continuous when its argument
passes through $1$ or tends to $\infty$, and its one-sided limits agree.
The remaining $\eta$ terms have at most $\log|t-1|$ singularities.
Integrating first on $(\varepsilon,1-\delta)$ and
$(1+\delta,R)$ and taking $\delta\downarrow0$ therefore produces no
additional interior boundary term. The endpoint limits computed above
and the convergent linear integrals justify subsequently taking
$\varepsilon\downarrow0$ and $R\to\infty$.
Combining the boundary and linear contributions yields
\[
 \int_0^\infty\mathcal T
 =-10D_\omega+\frac{11}{2}D_\omega
 =-\frac92D_\omega.
\]
Finally the absolutely convergent sine series gives
$\mathcal D(\omega)=\operatorname{Cl}_2(2\pi/3)=\sqrt3 K/2$.
Equation \eqref{s419:eq:sun419-S1-T} proves
\[
 \boxed{S_1=\frac{27}{4}K.}
\]
The independently established relation
$3S_1+S_2=4\pi\log3/\sqrt3-3K$ then gives
\[
 \boxed{S_2=\frac{4\pi\log3}{\sqrt3}-\frac{93}{4}K.}
\]
Thus both original identities of Sun's Conjecture~4.19 follow.

\clearpage
% Complete additional chapter; requires amsmath,amssymb,amsthm,hyperref.
\section{Both identities in Sun's Conjecture 4.27}
\label{sec:sun427}

The target here is the entire pair in Sun's Conjecture 4.27,
equations (4.28)--(4.29), in \emph{New series involving binomial
coefficients (II)}, arXiv:2307.03086v10. The proof uses a two-term
hypergeometric acceleration, followed by a residue at a logarithmically
divergent initial row. Two-term recursions of this kind were developed
by J. M. Campbell and P. Levrie, \emph{On two-term hypergeometric
recursions with free lower parameters}, J. Difference Equ. Appl.
30 (2024), 870--891, arXiv:2305.00626, and by J. M. Campbell,
\emph{Hypergeometric accelerations with shifted indices},
arXiv:2405.02776~\cite{CampbellLevrie,Campbell}. All recurrence identities, limiting operations, and
specializations needed below are proved explicitly. No historical
priority claim is made for the identities or the general method.

Put
\[
 t_j=\frac{\binom{3j}{j}\binom{4j}{j}^{\!2}}
 {(3j+1)(3j+2)4096^j},\qquad
 P(j)=368j^3+400j^2+118j+9.
\]
\begin{theorem}\label{thm:sun427-complete}
With $H_0=0$ and $H_r=\sum_{i=1}^{r}1/i$, both series converge
absolutely, and
\begin{align}
 \sum_{j=0}^{\infty}t_j\bigl\{P(j)(H_{2j}+H_j)
       -128j^2-136j-31\bigr\}
 &=-\frac{64\log2}{\pi},\label{eq:s427-target1}\\
 \sum_{j=0}^{\infty}t_j\bigl\{P(j)(2H_{4j}-5H_{2j})
       +192j^2+212j+46\bigr\}
 &=\frac{96\log2}{\pi}.\label{eq:s427-target2}
\end{align}
These are both original assertions in the numbered conjecture.
\end{theorem}

\subsection{An explicit acceleration certificate}

For parameters near $(a,b,n,m)=(1/2,1/2,1,1)$, define
\[
 \lambda=n+m-a-b-1,\qquad s=a+b,\qquad p=ab,
\]
and the following polynomials:
\begin{align*}
 T&=p+2s+1-n(m+1),\\
 D_1&=T+\lambda(n+2m+1),\\
 E&=s+p-n(m+1),\\
 D_2&=m(m+1)+n(2m+1),\\
 N&=\lambda(\lambda+1)D_2+D_1E+(s+2p)(\lambda+1).
\end{align*}
Use these definitions at each value of the displayed parameters, and set
\begin{align}
 U(a,b;n,m)&=\frac{N}{m(\lambda+1)(\lambda+2)},\label{eq:s427-U}\\
 v(a,b;n,m)&=
 \frac{(n-a)(n-b)(m-a)(m+1-a)(m-b)(m+1-b)}
 {nm(m+1)\lambda(\lambda+1)(\lambda+2)}.\label{eq:s427-v}
\end{align}
For the moment assume $\Re\lambda>0$, so that
\[
 f(a,b;n,m)=\sum_{k=0}^{\infty}
       \frac{(a)_k(b)_k}{(n)_k(m)_k}
\]
is absolutely convergent. Here $(z)_k=z(z+1)\cdots(z+k-1)$ and
$(z)_0=1$.

\begin{lemma}\label{lem:s427-recurrence}
For parameters sufficiently close to the indicated base point with
$\Re\lambda>0$,
\begin{equation}\label{eq:s427-recurrence}
 f(a,b;n,m)=\frac{U(a,b;n,m)}{\lambda}
       +v(a,b;n,m)f(a,b;n+1,m+2).
\end{equation}
\end{lemma}
\begin{proof}
Let $F_k=(a)_k(b)_k/((n)_k(m)_k)$ and write
\[
 X(k)=-\frac{k^2}{\lambda}
 -\frac{D_1 k}{\lambda(\lambda+1)}
 -\frac{N}{\lambda(\lambda+1)(\lambda+2)},\qquad
 R(k)=\frac{X(k)}{k+m}.
\]
The polynomial identity
\begin{align}
 &(k+a)(k+b)X(k+1)-(k+n)(k+m+1)X(k)\notag\\
 &\qquad=(k+n)(k+m)(k+m+1)-nm(m+1)v
 \label{eq:s427-polycert}
\end{align}
holds identically. For an explicit way to verify it, compare the
coefficients of $k^3,k^2,k$, which successively give the three
coefficients of $X$ above. The constant coefficient, after collecting
factors, is precisely \eqref{eq:s427-v}. Thus this is a rational
identity with unrestricted parameters, not a check at finitely many
integer arguments.

Since
\[
 \frac{F_{k+1}}{F_k}=\frac{(k+a)(k+b)}{(k+n)(k+m)},\qquad
 \frac{(a)_k(b)_k/((n+1)_k(m+2)_k)}{F_k}
 =\frac{nm(m+1)}{(k+n)(k+m)(k+m+1)},
\]
\eqref{eq:s427-polycert} gives
\[
 F_{k+1}R(k+1)-F_kR(k)
 =F_k-v\frac{(a)_k(b)_k}{(n+1)_k(m+2)_k}.
\]
The gamma-ratio asymptotic is $F_k=O(k^{-1-\Re\lambda})$,
whereas $R(k)=O(k)$. Consequently the upper telescoping boundary
vanishes when $\Re\lambda>0$. The lower boundary is
$-R(0)=U/\lambda$, proving the recurrence.
\end{proof}

\subsection{The balanced residue and its uniform justification}

Define an analytic coefficient, including at $\lambda=0$, by
\begin{equation}\label{eq:s427-W}
 W_j(a,b;n,m)=
 \frac{4^j(n-a)_j(n-b)_j
 (\frac{m-a}{2})_j(\frac{m+1-a}{2})_j
 (\frac{m-b}{2})_j(\frac{m+1-b}{2})_j}
 {(n)_j(\frac m2)_j(\frac{m+1}{2})_j(\lambda+1)_{3j}}.
\end{equation}
In particular $W_0=1$. None of its denominators vanish in a sufficiently
small complex neighborhood of the base point.

\begin{lemma}\label{lem:s427-residue}
If $(a,b,n,m)$ is sufficiently close to $(1/2,1/2,1,1)$ and
$n+m=a+b+1$, then
\begin{equation}\label{eq:s427-residue}
 \sum_{j=0}^{\infty}W_j(a,b;n,m)
             U(a,b;n+j,m+2j)
 =\frac{\Gamma(n)\Gamma(m)}{\Gamma(a)\Gamma(b)}.
\end{equation}
The series is locally uniformly convergent in these parameters, and
it may be differentiated along any analytic parameter curve lying
on the displayed balanced surface.
\end{lemma}
\begin{proof}
First take $\Re\lambda>0$. Direct multiplication of
\eqref{eq:s427-v} shows
\[
 W_j=\frac{\lambda}{\lambda+3j}
       \prod_{r=0}^{j-1}v(a,b;n+r,m+2r).
\]
This is also true for $j=0$, with the empty product interpreted as $1$.
Iterating \eqref{eq:s427-recurrence} a finite number of times gives
\begin{align}
 \lambda f(a,b;n,m)
 &=\sum_{j=0}^{J-1}W_jU(a,b;n+j,m+2j)\notag\\
 &\quad+(\lambda+3J)W_J f(a,b;n+J,m+2J).
 \label{eq:s427-finiteiter}
\end{align}

The exact rational ratio obtained from \eqref{eq:s427-W} satisfies
$W_{j+1}/W_j\longrightarrow4/27$ uniformly on a sufficiently small
closed complex neighborhood of the base point. All its numerator
factors stay nonzero there. Hence, for any fixed $q$ with
$4/27<q<1$, there is a uniform constant $C$ such that
$|W_j|\le Cq^j$. Formula \eqref{eq:s427-U} gives
$|U(a,b;n+j,m+2j)|\le C(j+1)$ on the same neighborhood.
Thus the accelerated series is holomorphic and locally uniformly
convergent, with the same property for its parameter derivatives.
For the remainder in \eqref{eq:s427-finiteiter}, the shifted sums
$f(a,b;n+J,m+2J)$ are uniformly bounded for $J\ge1$. For example,
after shrinking the neighborhood, their absolute summands are bounded
by a constant multiple of
$(3/4)_k^2/((7/4)_k(11/4)_k)$, a convergent series independent of $J$.
It follows that the remainder tends to zero.

It remains to evaluate the limit as $\lambda\to0$; the divergent
initial row is not discarded. Take real parameters on the balanced
surface near the base point, and replace $m$ by $m+\varepsilon$,
where $\varepsilon>0$. Uniform gamma-ratio asymptotics give, for $k\ge1$,
\[
 \frac{(a)_k(b)_k}{(n)_k(m+\varepsilon)_k}
 =C(\varepsilon)k^{-1-\varepsilon}
   \bigl(1+O(k^{-1})\bigr),\qquad
 C(\varepsilon)=\frac{\Gamma(n)\Gamma(m+\varepsilon)}
 {\Gamma(a)\Gamma(b)},
\]
where the error is uniform for small $\varepsilon\ge0$.
The elementary integral comparison
$\varepsilon\sum_{k=1}^{\infty}k^{-1-\varepsilon}\to1$
therefore yields
\[
 \lim_{\varepsilon\downarrow0}
 \varepsilon f(a,b;n,m+\varepsilon)
 =\frac{\Gamma(n)\Gamma(m)}{\Gamma(a)\Gamma(b)}.
\]
Uniform geometric convergence permits the same limit term by term
on the accelerated side of \eqref{eq:s427-finiteiter}, proving
\eqref{eq:s427-residue} for real parameters near the base point.
Holomorphy extends the identity to the complex balanced neighborhood;
alternatively its real parameter curves can be differentiated directly
using the local uniform derivative bounds just proved.
\end{proof}

\subsection{The two harmonic directions}

Write $D(j)=(3j+1)(3j+2)$. At the base point, elementary finite-product
cancellation and \eqref{eq:s427-U} give
\begin{equation}\label{eq:s427-base}
 W_j(1/2,1/2;1,1)=t_jD(j),\qquad
 U(1/2,1/2;1+j,1+2j)=\frac{P(j)}{16D(j)}.
\end{equation}
For the first identity use the balanced curve
\[
 (a,b,n,m)=(1/2-2z,\,1/2,\,1-z,\,1-z),
\]
and for the second use
\[
 (a,b,n,m)=(1/2+3z,\,1/2,\,1+z,\,1+2z).
\]
Both curves satisfy $n+m=a+b+1$ for every $z$ close to zero.
Denote differentiation along them at zero by $\partial_1$ and
$\partial_2$, respectively.

The finite-product identity
\[
 \frac{d}{dx}\log(x)_j=\sum_{r=0}^{j-1}\frac1{x+r}
\]
gives the exact formulas
\begin{align}
 \partial_1\log W_j&=H_{2j}+H_j,&
 \partial_2\log W_j&=2H_{4j}-5H_{2j}.
 \label{eq:s427-harmonic}
\end{align}
For transparency, set
\[
 A_j=\sum_{r=0}^{j-1}\frac1{r+1/2}=2H_{2j}-H_j,\qquad
 B_j=\sum_{r=0}^{j-1}\left(\frac1{r+1/4}+\frac1{r+3/4}\right)
       =4H_{4j}-2H_{2j}.
\]
The derivatives in \eqref{eq:s427-harmonic} reduce, respectively,
to $3H_j/2+A_j/2$ and $B_j/2-2A_j-2H_j$.
These formulas include $j=0$, where all the sums are empty.

Differentiating the explicit rational formula \eqref{eq:s427-U}, with
its shifted arguments $n+j,m+2j$, gives the polynomial identities
\begin{align}
 16D(j)\,\partial_1 U&=-128j^2-136j-31,
 \label{eq:s427-Q1}\\
 16D(j)\,\partial_2 U&=192j^2+212j+46.
 \label{eq:s427-Q2}
\end{align}
Equations \eqref{eq:s427-base}--\eqref{eq:s427-Q2} imply that sixteen
times the two derivatives of the left side of
\eqref{eq:s427-residue} are precisely the two series in
Theorem~\ref{thm:sun427-complete}.

The right side of \eqref{eq:s427-residue} equals $1/\pi$ at zero.
With $\psi=\Gamma'/\Gamma$, the standard gamma duplication formula
implies $\psi(1/2)=\psi(1)-2\log2$. Thus its two logarithmic
derivatives are
\[
 -2\psi(1)+2\psi(1/2)=-4\log2,
 \qquad
 3\psi(1)-3\psi(1/2)=6\log2.
\]
Multiplying the resulting derivatives by sixteen proves
\eqref{eq:s427-target1} and \eqref{eq:s427-target2}.
Absolute convergence also follows from the exact positive ratio
\[
 \frac{t_{j+1}}{t_j}
 =\frac{(4j+1)^2(2j+1)(4j+3)^2}
 {384(j+1)^3(3j+4)(3j+5)}\longrightarrow\frac4{27},
\]
since the remaining factors grow at most polynomially times
$1+\log(j+1)$.

\paragraph{Sources and verification scope.}
The original target is at
\url{https://arxiv.org/html/2307.03086v10}, Conjecture 4.27.
Primary methodological antecedents are
\url{https://arxiv.org/pdf/2305.00626} and
\url{https://arxiv.org/pdf/2405.02776}.
The certificate, base specialization, and both correction polynomials
have also been checked as exact rational identities by symbolic
simplification. The proof of the infinite identities is the finite
recurrence, the geometric remainder bound, the explicitly evaluated
residue, and justified differentiation; finite numerical agreement is
not used as a substitute for any of those steps.

\clearpage
% Complete original-target proof fragment. No external input files are required.
% Parent supplies amsmath, amssymb, amsthm, theorem/lemma environments, and \url.
\section{Both identities in Sun's Conjecture 4.25}
\label{sec:sun425-complete}

This section proves both identities in Conjecture 4.25 of Z.-W. Sun,
\emph{New series involving binomial coefficients (II)},
arXiv:2307.03086v10 (6 February 2026), equations (4.26)--(4.27).
The non-harmonic antecedent is attributed in Sun's Remark 4.26 to
Chu--Zhang, Example 110. The acceleration and parameter-differentiation
methods used here are classical. We supply the rational certificate and
all limiting arguments, and make no claim of historical priority.
The already known pair in Conjecture 4.13 is not part of this result.

Write $H_0=0$, $H_n=\sum_{r=1}^n1/r$, and put
\[
 t_j=\frac{\binom{3j}{j}\binom{4j}{j}^{2}}
 {(3j+1)(3j+2)1024^j},\qquad
 P(j)=88j^3+108j^2+36j+3.
\]
\begin{theorem}\label{thm:sun425-complete}
Both series below converge absolutely, and
\begin{align}
 \sum_{j=0}^{\infty}t_j\bigl\{
 P(j)(2H_{2j}-3H_j)+40j^2+56j+14\bigr\}
 &=\frac{32\log2}{\pi},\label{eq:s425-first}\\
 \sum_{j=0}^{\infty}t_j\bigl\{
 P(j)(2H_{4j}-3H_{2j})+40j^2+32j+8\bigr\}
 &=\frac{16\log2}{\pi}.\label{eq:s425-second}
\end{align}
These are the two complete assertions of the specified conjecture.
\end{theorem}

\subsection{A rational acceleration certificate}
For parameters $a,b,x,y$ define
\[
 \lambda=x+y-a-b-1,\qquad s=x+y,\qquad p=xy.
\]
Initially take real parameters in a sufficiently small neighbourhood of
$(a,b,x,y)=(1/2,1/2,1,1)$ with $\lambda>0$. In particular,
$0<a<1$, $b>0$, and $x,y>b$ can be assumed. Define
\[
 F_{a,b,x,y}(k)=\frac{(a)_k(b)_k}{(x)_k(y)_k},\qquad
 f(a,b,x,y)=\sum_{k=0}^{\infty}F_{a,b,x,y}(k),
\]
where $(z)_k=z(z+1)\cdots(z+k-1)$ and $(z)_0=1$.
The summand is $O(k^{-1-\lambda})$, so this series converges absolutely.

The following explicit expressions are convenient for keeping the
certificate short:
\begin{align*}
 E&=ab+2a+b-1-p+\lambda(s+a-1),\\
 J&=a-b-1+2ab,\qquad M=a-1+ab-p,\\
 \mathcal N&=(\lambda+1)J+EM
       +\lambda(\lambda+1)\{p+(a-1)s\},\\
 X(k)&=-\frac{k^2}{\lambda}
       -\frac{Ek}{\lambda(\lambda+1)}
       -\frac{\mathcal N}{\lambda(\lambda+1)(\lambda+2)},\\
 R(k)&=\frac{X(k)}{k+a-1},\\
 v&=\frac{(x-a)(x-a+1)(y-a)(y-a+1)(x-b)(y-b)}
 {(1-a)xy\lambda(\lambda+1)(\lambda+2)},\\
 U&=\frac{\mathcal N}{(a-1)(\lambda+1)(\lambda+2)},
 \qquad u=\frac{U}{\lambda}.
\end{align*}
All expressions denote rational functions of the displayed parameters.
The formulas for $U$ and $v\lambda$ are regular at $\lambda=0$
in the neighbourhood under consideration.

\begin{lemma}\label{lem:s425-recurrence}
For the above real parameters with $\lambda>0$,
\[
 f(a,b,x,y)=u+v\,f(a-1,b,x+1,y+1).
\]
The same recurrence holds at every shifted parameter tuple
$(a-N,b,x+N,y+N)$ for integers $N\ge0$.
\end{lemma}
\begin{proof}
The certificate is the following polynomial identity in $k$:
\begin{equation}\label{eq:s425-polycert}
\begin{split}
 &(k+a-1)(k+b)X(k+1)-(k+x)(k+y)X(k)\\
 &\hspace{12mm}=(k+a-1)(k+x)(k+y)-v(a-1)xy.
\end{split}
\end{equation}
It is verified by substituting the displayed expressions and comparing
the four coefficients of $1,k,k^2,k^3$; there are no numerical
or finite-range assumptions in this identity. In particular, the leading
coefficient gives $-\lambda(-1/\lambda)=1$, and the subsequent
coefficients give successively the displayed $E$, $\mathcal N$, and $v$.

The summand ratios are
\[
 \frac{F_{a,b,x,y}(k+1)}{F_{a,b,x,y}(k)}
  =\frac{(k+a)(k+b)}{(k+x)(k+y)},
 \quad
 \frac{F_{a-1,b,x+1,y+1}(k)}{F_{a,b,x,y}(k)}
  =\frac{(a-1)xy}{(k+a-1)(k+x)(k+y)}.
\]
Consequently \eqref{eq:s425-polycert} is precisely
\begin{equation}\label{eq:s425-telescope}
\begin{split}
 &F_{a,b,x,y}(k+1)R(k+1)-F_{a,b,x,y}(k)R(k)\\
 &\hspace{12mm}=F_{a,b,x,y}(k)
       -vF_{a-1,b,x+1,y+1}(k).
\end{split}
\end{equation}
Here $R(k)=O(k)$, so the term at infinity is
$O(k^{-\lambda})$ and vanishes. The term at zero is
$-R(0)=u$. Summation proves the recurrence.

Under the shift, $\lambda$ becomes $\lambda+3N>0$.
The parameter $a-N$ remains a noninteger near a half-integer,
so no denominator $k+a-N-1$ vanishes at a nonnegative integer $k$.
The same gamma-ratio estimate and rational identity therefore prove
all the asserted shifted recurrences.
\end{proof}

\subsection{Iteration, its remainder, and the balanced boundary}
Let
\[
 \theta_j=(a-j,b,x+j,y+j),\qquad
 U_j=U(\theta_j),\qquad v_j=v(\theta_j).
\]
Define $W_0=1$ and, initially for $\lambda>0$,
\[
 W_j=\frac{\lambda}{\lambda+3j}\prod_{r=0}^{j-1}v_r.
\]
The product has the following expression, which is also regular at
$\lambda=0$:
\begin{equation}\label{eq:s425-W}
 W_j=
 \frac{16^j
 (\frac{x-a}{2})_j(\frac{x-a+1}{2})_j
 (\frac{y-a}{2})_j(\frac{y-a+1}{2})_j
 (x-b)_j(y-b)_j}
 {(x)_j(y)_j(1-a)_j(\lambda+1)_{3j}}.
\end{equation}
This follows by grouping the four factors whose arguments increase by
$2$ and using
$\lambda(\lambda)_{3j}^{-1}/(\lambda+3j)
=(\lambda+1)_{3j}^{-1}$.
For every finite $N$, iteration of Lemma~\ref{lem:s425-recurrence} gives
\begin{equation}\label{eq:s425-iteration}
 \lambda f(a,b,x,y)
 =\sum_{j=0}^{N-1}W_jU_j
   +(\lambda+3N)W_N f(a-N,b,x+N,y+N).
\end{equation}

We now bound the last term, including the change of sign in the first
numerator parameter. For $N\ge1$ the two Euler beta integrals give
\begin{equation}\label{eq:s425-beta}
 f(a-N,b,x+N,y+N)=\mathbb E[(1-ST)^{N-a}],
\end{equation}
where $S,T$ are independent with beta distributions
\[
 S\sim\operatorname{Beta}(1,x+N-1),\qquad
 T\sim\operatorname{Beta}(b,y+N-b).
\]
All beta parameters are positive in the chosen real neighbourhood.
To verify the formula, expand $(1-ST)^{N-a}$ by the binomial series
and use
\[
 \mathbb E[S^k]=\frac{k!}{(x+N)_k},\qquad
 \mathbb E[T^k]=\frac{(b)_k}{(y+N)_k}.
\]
The sum of integrals of the absolute values of these terms is finite,
because its terms are the absolute values of the defining series with
excess $\lambda+3N>0$. This justifies the interchange, including at the
integrable corner $S=T=1$.
Since $N-a>0$ and $0\le ST\le1$, \eqref{eq:s425-beta} proves
\begin{equation}\label{eq:s425-fN}
 0<f(a-N,b,x+N,y+N)\le1.
\end{equation}

The ratio of consecutive terms in \eqref{eq:s425-W} tends to
$16/27<1$, uniformly when the initial parameters vary over a sufficiently
small closed neighbourhood of the base point. Thus, for some $\rho<1$
and constant $C$, $|W_j|\le C\rho^j$ there. The rational expression for
$U_j$ and each of its first parameter derivatives grows at most
polynomially in $j$. First derivatives of $W_j$ have at most an
additional $O(1+\log(j+1))$ factor, as follows by logarithmic
differentiation of its Pochhammer factors. Therefore both the accelerated
series and its first parameter derivatives converge uniformly.
Equations \eqref{eq:s425-fN} and \eqref{eq:s425-iteration} imply
\begin{equation}\label{eq:s425-accelerated}
 \lambda f(a,b,x,y)=\sum_{j=0}^{\infty}W_jU_j
 \qquad(\lambda>0).
\end{equation}
In particular, the iterated remainder has been proved to vanish;
it is not assumed from a formal transformation.

For completeness, the boundary at $\lambda=0$ is explicit. Uniformly
on a compact neighbourhood of the base point, the gamma-ratio
asymptotic gives, for $k\ge1$,
\[
 F_{a,b,x,y}(k)=
 \frac{\Gamma(x)\Gamma(y)}{\Gamma(a)\Gamma(b)}
 k^{-1-\lambda}\{1+O(k^{-1})\}.
\]
The error sum is uniformly bounded for real $\lambda\ge0$ near zero,
and the integral comparison for $\sum_{k\ge1}k^{-1-\lambda}$ gives
\[
 \lambda\sum_{k\ge1}k^{-1-\lambda}\longrightarrow1
 \qquad(\lambda\downarrow0).
\]
Multiplying the previous asymptotic by $\lambda$ and summing proves
\[
 \lambda f(a,b,x,y)\longrightarrow
 \frac{\Gamma(x)\Gamma(y)}{\Gamma(a)\Gamma(b)}.
\]
Taking the limit in the uniformly convergent right side of
\eqref{eq:s425-accelerated} consequently yields the identity
\begin{equation}\label{eq:s425-balanced}
 \boxed{\displaystyle
 \sum_{j=0}^{\infty}W_jU_j
 =\frac{\Gamma(x)\Gamma(y)}{\Gamma(a)\Gamma(b)}
 \quad\text{when }x+y-a-b-1=0,}
\end{equation}
for all real parameter tuples on that hyperplane sufficiently close
to $(1/2,1/2,1,1)$. The uniform first-derivative estimates above permit
differentiation along either of the tangent directions used next.

\subsection{The two independent harmonic directions}
At the base point, factorial cancellation in \eqref{eq:s425-W} and
substitution in $U$ give
\begin{equation}\label{eq:s425-base}
 W_j=(3j+1)(3j+2)t_j,
 \qquad
 U_j=\frac{P(j)}{8(3j+1)(3j+2)}.
\end{equation}
Thus the zeroth-order case of \eqref{eq:s425-balanced} also proves
$\sum_j t_jP(j)=8/\pi$.
Consider the tangent differential operators
\[
 D_1=\partial_a+\partial_b+\partial_x+\partial_y,
 \qquad D_2=\partial_b+\partial_x.
\]
Both annihilate $\lambda$. At the base point, logarithmic
differentiation of \eqref{eq:s425-W} gives
\begin{align}
 D_1W_j&=W_j(2H_{2j}-3H_j),\label{eq:s425-dW1}\\
 D_2W_j&=W_j(2H_{4j}-3H_{2j}).\label{eq:s425-dW2}
\end{align}
For an explicit check, write a tangent displacement as
$(\delta a,\delta b,\delta x,\delta y)=(A,B,X,Y)$ with
$X+Y=A+B$. The derivative of $\log W_j$ is
\[
 2(B-A)H_{4j}+(5A-3B)H_{2j}-3AH_j.
\]
This follows directly from
\[
 \sum_{r=0}^{j-1}\frac1{r+1/2}=2H_{2j}-H_j,
 \quad
 \sum_{r=0}^{j-1}\left(\frac1{r+1/4}+\frac1{r+3/4}\right)
 =4H_{4j}-2H_{2j}.
\]
Taking $(A,B)=(1,1)$ and $(A,B)=(0,1)$ gives the two displayed
formulas, including the empty sums at $j=0$.

The remaining derivatives involve only the displayed rational function
$U$. Substitution and clearing denominators give the exact identities
\begin{align}
 8(3j+1)(3j+2)D_1U_j&=40j^2+56j+14,
 \label{eq:s425-dU1}\\
 8(3j+1)(3j+2)D_2U_j&=40j^2+32j+8.
 \label{eq:s425-dU2}
\end{align}
All derivatives in \eqref{eq:s425-dW1}--\eqref{eq:s425-dU2} are
evaluated at the base point. Equations \eqref{eq:s425-base}--
\eqref{eq:s425-dU2} therefore show that the two left sides in
Theorem~\ref{thm:sun425-complete} are respectively
\[
 8D_1\sum_{j\ge0}W_jU_j,
 \qquad 8D_2\sum_{j\ge0}W_jU_j.
\]

Finally, set $C(a,b,x,y)=\Gamma(x)\Gamma(y)/(\Gamma(a)\Gamma(b))$.
At the base point $C=1/\pi$. If $\psi=\Gamma'/\Gamma$, the
logarithmic derivative of the gamma duplication formula gives
$\psi(1)-\psi(1/2)=2\log2$. Hence
\[
 D_1C=\frac{4\log2}{\pi},\qquad
 D_2C=\frac{2\log2}{\pi}.
\]
Differentiating \eqref{eq:s425-balanced} and multiplying by $8$
proves \eqref{eq:s425-first} and \eqref{eq:s425-second}.
Absolute convergence follows from the same geometric bounds already
proved. No numerical identification of a constant or unproved boundary
value enters this derivation.

\paragraph{Primary sources and verification scope.}
The exact original target is at
\url{https://arxiv.org/html/2307.03086v10}, Conjecture 4.25.
For the established acceleration background, see J.~M.~Campbell,
\emph{Hypergeometric accelerations with shifted indices},
\url{https://arxiv.org/pdf/2405.02776}, and the Chu--Zhang reference
identified in Sun's Remark 4.26. The rational certificate and the
rational derivative identities above were additionally checked by exact
symbolic algebra. Those checks confirm the printed formulas; the
convergence, remainder, and boundary arguments are supplied in full here.
\clearpage\appendix
\section{Coverage of the original twenty rounds}
The archive contains 120 participant submissions, six in each of twenty
rounds, plus a separate independent audit of the Showa Round 10 paper.
Six participant submissions had already been selected as complete in the
preceding audit. This pass revisited the remaining 114 submissions, joining
repeated rounds on a common target before testing the missing implication.
The table below records original-target completion in this pass, not the
truth or falsity of auxiliary lemmas in a manuscript.

\begin{center}
\begin{tabular}{@{}rcccccc@{}}\toprule
Round & Keio & Meiji & Taisho & Showa & Heisei & Reiwa\\\midrule
1 & $\cdot$ & $\cdot$ & $\cdot$ & $\cdot$ & $\cdot$ & $\cdot$\\
2 & $\cdot$ & $\cdot$ & $\cdot$ & $\cdot$ & $\cdot$ & $\cdot$\\
3 & $\cdot$ & $\cdot$ & $\cdot$ & $\cdot$ & $\cdot$ & P\\
4 & $\cdot$ & $\cdot$ & $\cdot$ & $\cdot$ & $\cdot$ & C\\
5 & $\cdot$ & $\cdot$ & $\cdot$ & $\cdot$ & $\cdot$ & C\\
6 & $\cdot$ & $\cdot$ & $\cdot$ & $\cdot$ & $\cdot$ & P\\
7 & $\cdot$ & $\cdot$ & $\cdot$ & $\cdot$ & $\cdot$ & P\\
8 & $\cdot$ & $\cdot$ & $\cdot$ & $\cdot$ & $\cdot$ & P\\
9 & $\cdot$ & $\cdot$ & $\cdot$ & $\cdot$ & $\cdot$ & C\\
10 & $\cdot$ & $\cdot$ & $\cdot$ & P & $\cdot$ & C\\
11 & $\cdot$ & $\cdot$ & $\cdot$ & $\cdot$ & E & P\\
12 & $\cdot$ & $\cdot$ & $\cdot$ & $\cdot$ & $\cdot$ & C\\
13 & $\cdot$ & $\cdot$ & $\cdot$ & $\cdot$ & $\cdot$ & $\cdot$\\
14 & $\cdot$ & $\cdot$ & $\cdot$ & $\cdot$ & $\cdot$ & $\cdot$\\
15 & $\cdot$ & $\cdot$ & $\cdot$ & $\cdot$ & $\cdot$ & $\cdot$\\
16 & $\cdot$ & $\cdot$ & $\cdot$ & $\cdot$ & $\cdot$ & $\cdot$\\
17 & $\cdot$ & $\cdot$ & $\cdot$ & $\cdot$ & $\cdot$ & $\cdot$\\
18 & $\cdot$ & $\cdot$ & $\cdot$ & $\cdot$ & $\cdot$ & $\cdot$\\
19 & $\cdot$ & $\cdot$ & $\cdot$ & $\cdot$ & $\cdot$ & $\cdot$\\
20 & $\cdot$ & $\cdot$ & $\cdot$ & $\cdot$ & $\cdot$ & $\cdot$\\
\bottomrule\end{tabular}
\end{center}
\noindent
\textbf{C}: a newly completed original target, proved in this volume.
\textbf{P}: previously completed submission, not counted again.
\textbf{E}: an external full-proof preprint was found; not counted as a
new proof produced here.
A dot means this pass did not produce a full proof of that submission's
original target. It does not assert that the problem is unresolved in every
external source, nor that its auxiliary results are incorrect.

Repeated number-theory rounds were tested against their remaining uniform
conditions: unbounded support, size-dependent congruences, simultaneous
positivity and height, and existence of a witness for every parameter.
Graph and matrix rounds were similarly tested at their remaining extension,
rank, or host-uniform steps. Finite computations and obstructions to a
restricted strategy were not promoted to original-target completions.

For Heisei Round 11, the current source search found Teng Zhang's
\emph{A proof of Schmeisser's conjecture}, arXiv:2611.19084v2, revised
17 November 2026~\cite{Zhang}. It claims the full convex-hull critical-point
statement in every degree, using an analytic argument and a finite
computer-assisted verification. This volume does not reproduce or certify
that external computational proof, and it is not included in our new-result
count. No theorem in the preceding chapters depends on it.

\paragraph{Verification scope.}
The displayed algebraic certificates were checked exactly, and the
improper endpoints, branch choices, and infinite-sum estimates were reviewed
separately from the symbolic calculations. The complete source compiles
without external bibliography files, figures, or auxiliary proof inputs.
Internal independent review is not represented as external peer review or
proof-assistant formalization.

\begin{thebibliography}{9}\raggedright
\bibitem{Sun}
Z.-W. Sun, \emph{New series involving binomial coefficients (II)},
arXiv:2307.03086v10, revised 6 February 2026.
Exact original statements are in Section 4.
\url{https://arxiv.org/html/2307.03086v10}.

\bibitem{Guillera}
J. Guillera, \emph{Hypergeometric identities for 10 extended
Ramanujan-type series}, Ramanujan Journal 15 (2008), 219--234.
Author preprint: arXiv:1104.0396 (2011).
\url{https://arxiv.org/abs/1104.0396}.

\bibitem{CampbellLevrie}
J. M. Campbell and P. Levrie,
\emph{On two-term hypergeometric recursions with free lower parameters},
Journal of Difference Equations and Applications 30 (2024), 870--891.
\url{https://arxiv.org/pdf/2305.00626}.

\bibitem{Campbell}
J. M. Campbell, \emph{Hypergeometric accelerations with shifted indices},
arXiv:2405.02776 (2024).
\url{https://arxiv.org/pdf/2405.02776}.

\bibitem{Zhang}
T. Zhang, \emph{A proof of Schmeisser's conjecture},
arXiv:2611.19084v2, 17 November 2026.
External proof claim; not a premise of this volume.
\url{https://arxiv.org/abs/2611.19084v2}.
\end{thebibliography}
\end{document}
